\documentclass[11pt]{article}

\usepackage[margin=1in]{geometry}
\usepackage{amsmath,amssymb,amsthm}

\usepackage{mathtools}
\usepackage{bm}
\usepackage{authblk}
\usepackage{enumitem}
\usepackage{booktabs}
\usepackage{graphicx}
\usepackage{microtype}
\usepackage[hidelinks]{hyperref}
\newcommand{\Ical}{\mathcal{I}}
\newcommand{\Dcal}{\mathcal{D}}
\newcommand{\Rcal}{\mathcal{R}}
\DeclareMathOperator{\Cert}{Cert}
\usepackage{xcolor}

% ---------- Macros ----------
\newcommand{\R}{\mathbb{R}}
\newcommand{\X}{\mathcal{X}}
\newcommand{\M}{\mathcal{M}}
\newcommand{\Pcal}{\mathcal{P}}
\newcommand{\Acal}{\mathcal{A}}
\newcommand{\E}{\mathbb{E}}
\newcommand{\Var}{\mathrm{Var}}
\newcommand{\Prb}{\mathbb{P}}
\newcommand{\Hit}{\operatorname{Hit}}
\newcommand{\dg}{d_g}
\newcommand{\deucl}{\delta}
\newcommand{\Ycal}{\mathcal{Y}}
\newcommand{\Ccal}{\mathcal{C}}
\newcommand{\Hcal}{\mathcal{H}}

\newcommand{\mirador}{\textsc{Mirador}}
\newcommand{\herald}{\textsc{Herald}}
\newcommand{\geodesic}{\textsc{Geodesic}}
\newcommand{\tessera}{\textsc{Tessera}}
\newcommand{\davis}{\textsc{Davis}}

\newtheorem{definition}{Definition}
\newtheorem{theorem}{Theorem}
\newtheorem{proposition}{Proposition}
\newtheorem{lemma}{Lemma}
\newtheorem{assumption}{Assumption}

% ---------- Title & Author ----------
\title{MIRADOR: Manifold of Immune Receptors And Drift-Optimized Responses\\[4pt]
Geometry-First Co-Embedding for Vaccine and Antibody Design}

\author{Bee Rosa Davis\thanks{%
  Correspondence: \texttt{bee\_davis@alumni.brown.edu}\quad
  ORCID: \texttt{0009-0009-8034-4308}.\\[2pt]
  \textbf{Patent notice.}
  Subject matter described in this manuscript may relate to
  US Provisional Patent Application No.\ 64/012,328:
  ``MIRADOR: Manifold-Informed Rational Architecture for Drug-Organism Response,
  Geometric Optimization of Therapeutic Coherence with Resistance Prediction via
  Escape Geodesics on Drug-Target Fiber Bundles.''
  Commercial use of the methods described herein may require a licence.
}}

\date{\today}

\begin{document}

\maketitle

\begin{abstract}
Many high-stakes problems in immunology involve \emph{two} coupled geometries: an antigenic
landscape traced out by pathogen evolution, and a complementary landscape of immune receptors
(antibodies, BCRs, TCRs) that must ``meet'' future antigens with sufficient binding energy.
Prior work in the \davis{} framework models identity-preserving processes as paths on a learned
Riemannian manifold with bounded geodesic--Euclidean distortion, abstention bands, and
compositional error budgets for \emph{detection} and \emph{surveillance}.\footnote{See
Davis manifolds for geometry-first detection and HERALD for viral surveillance. Citations are
suppressed in this draft to keep notation uncluttered.}

This paper introduces \mirador{}---the \emph{Manifold of Immune Receptors And
Drift-Optimized Responses}---as the design-side counterpart to geometry-first surveillance.
If \herald{} is a radar for antigenic drift and \geodesic{} is a filter for cancer signals, then
\mirador{} is an engineer: given a forecasted family of antigenic paths, it searches a
co-embedded receptor manifold for panels of antibodies or vaccine payloads that intercept
those paths with provable coverage guarantees.

Formally, we (i) define \emph{mirror manifolds}: coupled antigen and receptor manifolds
linked by a complementary co-embedding in which distance approximates \emph{binding free
energy}, not sequence similarity; (ii) construct contrastive, asymmetrically trained
co-embeddings that distinguish useless from mismatched receptors and implement a
lock-and-key geometry via maps $f_A$ and $f_R$ into a shared antigenic space; (iii) introduce
\emph{inverted error budgets} that bound \emph{minimum efficacy} (coverage of drift paths)
rather than maximum failure probability, including a vacuity floor where uncertainty forces
success guarantees back to zero; (iv) cast computational vaccine design as a set cover
problem on a Riemannian manifold, selecting receptor or payload panels that cover predicted
path families with certificates of coverage; and (v) incorporate a dual-use ``firewall''
that constrains \mirador{} to broad-spectrum panels, preventing the system from implicitly
leaking optimal sequences for unobserved, potentially malicious antigens.

Scope: This manuscript is conceptual. It provides definitions, assumptions, theorems, and
proposed evaluation protocols but does not report empirical performance, neutralization
results, or clinical outcomes. Within this scope, \mirador{} extends the \davis{}
blueprint from passive surveillance to active immune engineering, replacing ``bounded risk
of failure'' with ``certified lower bounds on success'' wherever the geometry is validated
and abstaining elsewhere.
\end{abstract}
\clearpage
\section{Introduction}

Modern vaccine design still treats antigenic evolution and immune response as loosely coupled
problems: we forecast how a virus will move in an antigenic landscape, and separately search for
antibodies or vaccine payloads that appear promising. This separation is increasingly untenable.
In systems where viral evolution is gradual and constrained, and where immune receptors adapt
via affinity maturation, both processes live on structured manifolds and interact through geometry.

Davis manifolds formalize this geometry-first view for detection problems.\footnote{See Davis
et al.\ for definitions of path families $\Pcal(L)$, distortion profiles $\varepsilon(L)$, and compositional
error budgets.} A Davis manifold is a Riemannian state space equipped with benign path families,
bounded geodesic--Euclidean distortion along those paths, and configuration margins tied to
abstention and error budgets. HERALD instantiates this framework for viral
surveillance: it constructs an antigenic pullback manifold on sequence space and uses path-wise
drift statistics, Cantelli bounds, and a compositional error budget to bound dominance risk.
In this sense, HERALD is a \emph{detector}: it forecasts where the pathogen is going and raises
alerts when drift crosses antigenic margins.

This paper takes the next step: from geometry-first detection to geometry-first \emph{design}.
We introduce \textbf{MIRADOR} (\emph{Manifold of Immune Receptors And Drift-Optimized
Responses}) as a mirror-system to HERALD. If HERALD is the radar that tracks the target,
MIRADOR is the engineer that positions the interceptors. Instead of bounding \emph{risk of failure}
for a surveillance statistic, MIRADOR aims to bound \emph{probability of success} for candidate
vaccine panels, under explicit assumptions and with an inverted error budget that collapses to
zero when geometric or modeling assumptions fail.

At the core of MIRADOR is a \emph{mirror} construction: a pair of coupled manifolds describing
antigenic drift and immune receptors, co-embedded into a shared ambient space so that distances
reflect \emph{interaction potential} rather than raw sequence similarity. We refer to this as a system
of \emph{mirror manifolds} only in a metaphorical sense---complementary antigen--receptor geometry
in the Davis style, not string-theoretic mirror symmetry.

\paragraph{Key idea.}
MIRADOR couples three elements:
(i) a complementary co-embedding of antigens and immune receptors in a shared latent space,
(ii) a path-wise notion of \emph{interception} for forecast antigenic drift paths by finite receptor
panels, and (iii) an \emph{inverted} Davis-style error budget that lower-bounds coverage of drift paths
by vaccine panels, with a hard floor at zero when the geometry or linkage becomes vacuous.
Formally, MIRADOR views vaccine design as a \emph{set-cover problem on a Riemannian
manifold}: choose a small set of receptors whose binding neighborhoods cover as much as possible
of a forecast antigenic path family, subject to geometric distortion and modeling uncertainty.

\subsection{From geometry-first detection to geometry-first design}

Davis manifolds and HERALD treat geometry as the substrate on which surveillance statistics
are built.\   Construction (contrastive training plus smoothness regularization)
provides a manifold $(\M, g)$ with path-class-specific distortion bounds $\varepsilon(L)$; traversal uses
Euclidean surrogates for geodesics inside that bounded-distortion regime; finite-variance Cantelli
bounds and compositional error budgets turn statistic-level separation into conditional risk
guarantees.

MIRADOR adopts the same construction–traversal–guarantee pipeline, but with a different
objective. Instead of asking
\[
\text{``How likely is it that a drift path has crossed a configuration margin?''}
\]
we ask
\[
\text{``Given a forecast family of drift paths and a vaccine panel $V$, how likely is it that
every such path is intercepted by $V$ within a binding radius?''}
\]
Here, interception is a path-wise geometric event (a drift trajectory passing within a binding
radius of some panel element), and the object of interest is a \emph{coverage} functional $C(V)$
over path families rather than a misclassification probability for a single observation.

\subsection{Mirror manifolds and complementary co-embedding}

We model antigens and immune receptors as distinct but coupled geometric objects.

\paragraph{State spaces and encoders.}
Let $\Acal$ denote the space of antigen configurations (e.g., viral spike sequences and structures),
and $\Rcal$ the space of receptor configurations (e.g., antibody or B-cell receptor sequences with
associated paratope features). We assume encoders
\[
f_A : \Acal \to \mathbb{R}^d, \qquad
f_R : \Rcal \to \mathbb{R}^d,
\]
whose images are endowed with pullback Riemannian metrics $g_A, g_R$ via the Euclidean metric
on $\mathbb{R}^d$. The resulting manifolds $(\M_A, g_A)$ and $(\M_R, g_R)$ inherit the Davis
structure: path families $\Pcal_A(L), \Pcal_R(L)$ of ``benign'' identity-preserving trajectories,
and path-class-specific distortion profiles $\varepsilon_A(L), \varepsilon_R(L)$ that bound geodesic–
Euclidean distortion along those paths.

Antigenic paths $\Gamma : [0,1] \to \M_A$ represent forecast drift trajectories with Riemannian
arc length
\[
\ell_{g_A}(\Gamma)
\;=\;
\int_0^1 \bigl\|\dot{\Gamma}(t)\bigr\|_{g_A}\, dt,
\]
so paths are \emph{not} assumed to have unit speed; the horizon constraint $\ell_{g_A}(\Gamma)\le L$
simply bounds intrinsic path length under $g_A$.

We write
\[
u_A(a) := f_A(a) \in \mathbb{R}^d, \qquad
u_R(r) := f_R(r) \in \mathbb{R}^d,
\]
and use the ambient Euclidean distance
\[
\delta\bigl(u_A(a), u_R(r)\bigr)
:= \bigl\|u_A(a) - u_R(r)\bigr\|_2
\]
as a proxy, within a bounded-distortion regime, for geodesic separation between antigen and
receptor manifolds.

\paragraph{Complementary rather than similarity embeddings.}
Naively co-embedding antigens and receptors and interpreting $\delta(u_A(a), u_R(r))$ as a
binding score risks conflating complementarity with similarity and admits trivial solutions where
all receptors collapse to a single latent centroid. MIRADOR therefore interprets $f_R$ as an
\emph{approximate map to the cognate antigen region}: $u_R(r)$ is trained to approximate the
location, in antigen space, of the ``platonic'' epitope that $r$ would ideally bind. Symmetrically,
$f_A$ places antigens so that known cognate receptors lie near their own inferred targets.

Concretely, the complementary structure is enforced via asymmetric contrastive training with
three classes of negatives:
(i) non-cognate antigens for a given receptor,
(ii) antigens that bind \emph{other} receptors (mismatched negatives), and
(iii) antigens that bind no receptor in the training panel (forcing specificity).
For each receptor $r$, we push $u_R(r)$ toward its known binding partners in antigen space and
away from all three negative types. This prevents trivial centroid collapse and encourages $u_R(r)$
to predict where its cognate antigens live in the antigen manifold (details in
Section~\ref{sec:mirror-construction}).

\paragraph{From distance to binding and neutralization.}
Within a calibrated regime, we posit a monotone surrogate map from distance to binding free
energy,
\[
\Phi_{\text{bind}}\!\bigl(\delta(u_A(a), u_R(r))\bigr)
\;\approx\;
-\Delta G_{\text{bind}}(a,r),
\]
where $\Delta G_{\text{bind}}(a,r) = G_{\text{complex}} - G_{\text{free}}$ is the Gibbs free-energy change on
binding, so more negative $\Delta G_{\text{bind}}$ means stronger binding. We take
$\Phi_{\text{bind}}$ to be decreasing so that smaller distances correspond to more favorable (more
negative) binding free energies.

We then distinguish \emph{binding} from \emph{neutralization} via a separate map
$\Phi_{\text{neut}}$ from predicted binding to neutralization probability (e.g., via epitope-level and
Fc-effector features),
\[
p_{\text{neut}}(a,r)
\;=\;
\Phi_{\text{neut}}\!\Bigl(\Phi_{\text{bind}}\!\bigl(\delta(u_A(a), u_R(r))\bigr),\; \text{context}\Bigr),
\]
and emphasize that MIRADOR's formal guarantees apply to neutralization only insofar as
$\Phi_{\text{neut}}$ is well-calibrated; residual mismatch is absorbed into a linkage error term
$E_{\text{link}}$ in the error budget.

A binding-distance threshold is defined by inverting $\Phi_{\text{bind}}$ at a pre-registered free-energy
cutoff:
\[
\rho_{\text{bind}}
\;:=\;
\Phi_{\text{bind}}^{-1}\!\bigl(\Delta G_{\text{thresh}}\bigr),
\]
where, for example, $\Delta G_{\text{thresh}} = -10 \text{ kcal/mol}$ might correspond to a strong-binding
regime under the chosen calibration.

\subsection{Interception, coverage, and set cover on a manifold}

We formalize vaccine panels as finite receptor sets and use path-wise interception as the basic
coverage event.

\paragraph{Forecast drift and interception.}
Let $\Pi_A(L)$ denote a probability distribution over benign antigenic drift paths
$\Gamma : [0,1]\to \M_A$ with $\ell_{g_A}(\Gamma)\le L$. We think of $\Pi_A(L)$ as a Davis-style
forecast family:\ it is a law over antigenic trajectories that respect biological plausibility constraints
(structural viability, fitness thresholds, etc.) to be formalized in Section~\ref{sec:path-families}.

A finite \emph{receptor panel} is a set $V \subset \Rcal$. Its image
$U_R(V) := \{u_R(r) : r \in V\}$ defines a family of binding neighborhoods in $\mathbb{R}^d$,
each of radius $\rho_{\text{bind}}$. Given a path $\Gamma$ with time horizon $T_\Gamma$ (e.g., a fixed
calendar horizon or implicit via arc-length parametrization), we say that $V$ \emph{intercepts}
$\Gamma$ if at some time $t$ the drift path enters one of these neighborhoods:
\[
\Hit(V,\Gamma)
:= \left\{\exists\, t \in [0, T_\Gamma],\ \exists\, r \in V :
\delta\bigl(u_A(\Gamma(t)), u_R(r)\bigr) \le \rho_{\text{bind}}\right\}.
\]

\paragraph{True coverage vs.\ empirical coverage.}
The \emph{true coverage} of a panel $V$ under the forecast family is the scalar
\[
C(V)
\;:=\;
\Pr_{\Gamma \sim \Pi_A(L)}\!\bigl(\Hit(V,\Gamma)\bigr),
\]
which is a deterministic function of $(V,\Pi_A(L))$ once both are fixed. In practice we can only
access an empirical estimate $\widehat{C}(V)$ constructed from $N$ sampled paths
$\{\Gamma_i\}_{i=1}^N \sim \Pi_A(L)$ or from a finite discretization of their state space. A MIRADOR
\emph{coverage certificate} with target level $C_0$ and confidence $1-\delta$ asserts that, with
probability at least $1-\delta$ over the sampling and residual model uncertainty,
\[
C(V) \;\ge\; C_0,
\]
where the gap between the nominal $C_0$ and the true $C(V)$ is controlled by an explicit lower
bound $L_{\text{MIR}}(V)$ (Section~\ref{sec:inverted-error-budgets}).

\paragraph{Set cover on a Riemannian manifold.}
To connect algebraically with classical set cover, MIRADOR works with a finite approximation to
the continuous drift family. We draw a finite set of representative antigenic states
$\{a_i\}_{i=1}^N$ along drift paths sampled from $\Pi_A(L)$, with sampling density chosen to ensure
that $|\widehat{C}(V) - C(V)| \le \varepsilon_{\text{samp}}$ with high probability (sampling schemes and
bounds are detailed in Section~\ref{sec:applications}). For each receptor $r$, define the discrete
coverage set
\[
S_r
:= \left\{ i \in \{1,\dots,N\} :
\delta\bigl(u_A(a_i), u_R(r)\bigr) \le \rho_{\text{bind}}\right\}.
\]
For a panel $V$, the union $\bigcup_{r\in V} S_r$ is the subset of sampled antigen states intercepted
by $V$, and the empirical coverage is
\[
\widehat{C}(V)
\;:=\;
\frac{1}{N}\, \bigl|\bigcup_{r\in V} S_r\bigr|.
\]

This is exactly a (weighted) set cover problem on a finite ground set:
given a budget $k$, choose a subset $V$ of receptors with $|V|\le k$ to maximize $\widehat{C}(V)$,
or, dually, to achieve a target coverage $\widehat{C}(V)\ge \widehat{C}_0$ with minimal $|V|$.
MIRADOR uses standard greedy algorithms and their classical approximation guarantees in this
finite setting, then analyzes how those guarantees transfer back to the continuous coverage $C(V)$
once geometric distortion and sampling error are folded into the error budget.

\subsection{Inverted error budgets and coverage certificates}

HERALD uses a \emph{forward} error budget: when geometric or linkage assumptions become
vacuous, the system becomes conservative by treating uncertainty as risk (raising or widening
alerts). MIRADOR must behave differently. When geometry or linkage is vacuous,
we cannot safely claim that a vaccine panel is good; we must instead declare that we cannot certify
coverage.

\paragraph{From statistic separation to coverage bounds.}
As in Davis systems, MIRADOR constructs scalar statistics (here, functions of
$\widehat{C}(V)$, binding scores, and linkage features) and uses finite-variance Cantelli bounds
to map statistic-level separation into probability statements. In particular, we define a Cantelli
term $B_{\text{Cantelli}}(V)$ for a coverage-related statistic, with the usual finite-variance form
(e.g.,
$B_{\text{Cantelli}} = 1 - \sigma^2/(\sigma^2 + \Delta^2)$ for variance $\sigma^2$ and separation $\Delta$),
as in the Davis analysis.\ 

\paragraph{Inverted Davis-style error budget.}
We then assemble an inverted error budget that decomposes failures into:
(i) geometry error $E_{\text{geom}}$ (distortion or path-family violations that break distance–binding
linkage),
(ii) linkage error $E_{\text{link}}$ (binding and neutralization modeling failures given good geometry),
(iii) calibration error $\xi$ (mismatch between reported and true coverage probabilities), and
(iv) abstention failure $\zeta$ (cases where the system should have refused to certify but did not),
together with an independence slack $\delta_{\text{indep}}$ to account for correlations between these
error sources, mirroring the Davis and HERALD budgets.\ 

For a panel $V$ and target level $C_0$, a MIRADOR coverage certificate is then backed by a
lower bound of the form
\[
\Pr\!\bigl(C(V) \ge C_0 \;\big|\; \Cert(V, C_0, 1-\delta)\bigr)
\;\ge\;
L_{\text{MIR}}(V),
\]
with
\[
L_{\text{MIR}}(V)
:= \max\Bigl\{0,\;
B_{\text{Cantelli}}(V)\,
\bigl(1-E_{\text{geom}}\bigr)\bigl(1-E_{\text{link}}\bigr)\bigl(1-\xi\bigr)\bigl(1-\zeta\bigr)
- \delta_{\text{indep}} - \delta\Bigr\}.
\]
Here the outer probability is taken over
(i) estimation noise in $\widehat{C}(V)$ from finite path sampling, and
(ii) residual epistemic uncertainty in the binding and neutralization models.
The $\max\{0,\cdot\}$ truncation enforces the desired floor: when geometry or linkage error becomes
large, or when the Cantelli term collapses, $L_{\text{MIR}}(V)$ goes to zero and no certificate is
issued. MIRADOR thus offers \emph{certificates of coverage}, not warnings: when assumptions
fail, the system refuses to certify rather than extrapolating optimistically.

\subsection{Scope, limitations, and dual-use firewall}

\paragraph{Scope and modeling approximations.}
This manuscript is conceptual. It defines mirror manifolds, complementary co-embeddings,
interception events, coverage functionals, and inverted error budgets; it states theorems and
proposes empirical protocols, but reports no empirical vaccine-design performance. Several
biological complexities are treated as modeling approximations whose errors are absorbed into
$E_{\text{link}}$ rather than as guaranteed features of the geometry:
(i) multi-epitope antigens and multi-paratope binding are handled via pooling or mixture models
over epitope-level distances,
(ii) conformational dynamics are represented by single conformers or small ensembles in latent
space, and
(iii) affinity maturation paths in receptor space are modeled via benign path families
$\Pcal_R(L)$ rather than fully stochastic dynamics. All guarantees are conditional on the
assumptions and audits specified later in the paper.

\paragraph{Dual-use and forbidden queries.}
A system that ties antigenic drift forecasts to receptor design is inherently dual-use: the same
geometry that helps design broad-coverage vaccines could, in principle, be inverted to search for
antigens that evade a given panel. MIRADOR is therefore designed with a dual-use firewall:

\begin{itemize}
  \item \emph{Input-tied optimization only.} Optimization is restricted to receptor panels
  $V$ and to antigenic states and path families explicitly provided as inputs. MIRADOR does
  not search over unconstrained antigen sequence space under abstract fitness objectives.

  \item \emph{No generative forecasting of bioweapons.} MIRADOR does not synthesize or rank
  hypothetical antigens that are not already specified by the user (e.g., as part of a forecast family
  $\Pi_A(L)$) and does not output stepwise mutation paths or mutation ``recipes.''

  \item \emph{No inverse or evasion queries.} The system refuses inverse queries of the form
  ``given panel $V$, find antigens that evade it'' or ``optimize an antigen sequence to minimize
  coverage by $V$.'' Even when a forward model exists that maps $(V,\Pi_A(L))$ to a coverage
  estimate, MIRADOR does not expose interfaces that allow adversarial inversion to recover
  high-risk evading sequences.

  \item \emph{Coarse-grained reporting.} Public or shared outputs are restricted to panel-level
  coverage summaries and coarse geometric attributions (e.g., coverage of drift families or regions),
  not fine-grained sequence-level recipes or unobserved combinatorial mutation sets.
\end{itemize}

Within these constraints, MIRADOR is intended as a theoretical capstone for the Davis program:
Davis manifolds supply the geometry and error budgets; HERALD instantiates
geometry-first detection for antigenic drift; MIRADOR adds a complementary,
mirror geometry for immune receptors and extends the Davis framework from bounding risk of
failure to certifying probability of success for vaccine panels, subject to explicitly stated and
auditable assumptions.

\paragraph{Contributions.}
Within this Davis-style perspective, the technical contributions of MIRADOR are:
\begin{itemize}
  \item \emph{Complementary mirror manifolds for antigen--receptor co-embedding.}
  We define coupled Davis manifolds $(\M_A, g_A)$ and $(\M_R, g_R)$ with encoders
  $f_A, f_R$ trained in an asymmetric contrastive fashion so that receptor embeddings predict
  locations of cognate antigens, and we formalize benign path families $\Pcal_A(L), \Pcal_R(L)$
  and distortion profiles $\varepsilon_A(L), \varepsilon_R(L)$ for this co-embedded setting.

  \item \emph{Interception as path-wise coverage.}
  We formalize interception events $\Hit(V,\Gamma)$ and coverage functionals $C(V)$ over
  forecast drift families, and show how finite set-cover approximations on sampled path
  skeletons connect classical greedy set cover guarantees back to continuous coverage
  under bounded distortion and sampling error.

  \item \emph{Inverted Davis-style error budgets and coverage certificates.}
  We derive MIRADOR's inverted error budget and associated lower bounds
  $L_{\text{MIR}}(V)$ for coverage certificates, combining finite-variance Cantelli terms with
  Davis-style error components $(E_{\text{geom}}, E_{\text{link}}, \xi, \zeta, \delta_{\text{indep}})$ in a way
  that collapses to zero when geometry or linkage is vacuous.

  \item \emph{Safety and dual-use constraints for geometric design.}
  We articulate input-tied optimization rules, forbidden inverse query classes, and coarse-grained
  reporting policies that allow MIRADOR to be analyzed as a geometric design framework
  without supplying mutation ``recipes'' or antigen-escape optimizers.
\end{itemize}

The remainder of the paper formalizes these elements: Section~\ref{sec:mirror-construction}
defines mirror manifolds and complementary co-embedding; Section~\ref{sec:path-families} formalizes
benign path families and distortion profiles in the co-embedded setting; Section~\ref{sec:inverted-error-budgets}
derives the inverted error budget and coverage certificates; and
Section~\ref{sec:applications} sketches how MIRADOR can be instantiated for geometric vaccine
design problems while respecting the dual-use firewall.

\subsection{Roadmap}

The rest of the paper is organized as follows.

\begin{itemize}[leftmargin=*]
  \item \textbf{Section 2 (Conceptual overview).} Introduces the mirror-manifold picture
  more formally, including the distinction between target and interceptor manifolds, the
  co-embedding maps $(f_A, f_R)$, and the notion of binding-energy geodesics.
  \item \textbf{Section 3 (Mirror manifolds: co-embedding).} Defines the mirror manifold,
  complementary maps, and the binding-energy metric. We show how to construct such
  manifolds from contrastive binding data with asymmetric negatives and smoothness
  regularization, yielding bounded-distortion regimes analogous to \davis{} manifolds.
  \item \textbf{Section 4 (Construction: algorithms and training objectives).} Specifies
  practical training objectives, including receptor/antigen sampling schemes, the treatment
  of useless vs.\ mismatched receptors, and regularization strategies that control the
  curvature and complementarity of the mirror geometry.
  \item \textbf{Section 5 (Inverted error budgets and certificates of coverage).} Develops
  the inverted Davis-style error budget, defines coverage certificates for receptor panels,
  and derives finite-variance lower bounds on success that decompose geometry, linkage,
  calibration, and abstention errors.
  \item \textbf{Section 6 (Applications and toy scenarios).} Illustrates how \mirador{}
  could be instantiated for computational vaccine design, including a toy ``universal flu''
  scenario framed as a center-of-mass problem on a drift path family, and a set-cover
  formulation for panel selection.
  \item \textbf{Section 7 (Safety, dual-use, and governance).} Describes the dual-use
  firewall in the design loop, including restrictions to broad-spectrum panels, coarse
  reporting of attributions, and limits on optimization over unobserved antigens.
  \item \textbf{Section 8 (Discussion and outlook).} Summarizes the role of \mirador{}
  within the broader \davis{} ecosystem, highlights limitations, and sketches open problems
  in joint antigen--receptor geometry, multi-modal co-embedding, and real-world validation.
\end{itemize}

Readers primarily interested in the high-level idea can focus on Sections 2 and 5;
those interested in the geometry and error budgets may wish to read Sections 3 and 5
in tandem; safety and governance aspects are collected in Section 7.

\section{Conceptual Overview and Mirror Geometry}
\label{sec:conceptual-overview}

MIRADOR lives at the ``design'' end of the Davis framework. Where HERALD treats the antigenic manifold as a space to \emph{forecast} viral drift and raise early-warning alerts, MIRADOR couples that geometry to a \emph{receptor} manifold and asks a different question:

\begin{quote}
  Given a forecast family of benign antigenic drift paths, where must a panel of immune receptors live in a shared geometric space to intercept those paths with certifiable probability?
\end{quote}

This section gives a coordinate-level overview of that coupling. We first position MIRADOR within the broader Davis suite, then specialize the Davis construction to antigen and receptor manifolds, and finally formalize interception as path-wise coverage in a complementary co-embedding.

\subsection{From forecasting to interception}

At a high level, the Davis suite can be viewed as a progression:

\begin{itemize}
  \item \textbf{Davis (geometry-first detection).} Provides a general recipe for constructing Riemannian manifolds with path-class-specific distortion bounds, configuration margins, and compositional error budgets for temporal detection tasks.

  \item \textbf{HERALD (radar).} Instantiates a Davis system on viral antigenic drift: a pullback manifold on sequence space, benign mutational paths, and a drift statistic whose separation admits finite-variance escape bounds and a compositional dominance error budget.

  \item \textbf{GEODESIC (filter).} Applies similar ideas to multi-cancer early detection, using geometry to denoise and interpret complex biological signals such as liquid-biopsy readouts.

  \item \textbf{TESSERA (network).} Extends the framework to microbial consortia and supply-chain style networks, modeling how perturbations propagate across coupled latent geometries.

  \item \textbf{MIRADOR (engineer).} Takes the antigen geometry as given and adds a coupled receptor geometry, asking \emph{where} receptors must be placed so that forecasted antigen paths are intercepted with certifiable coverage.
\end{itemize}

In HERALD, the latent geometry is purely \emph{observational}: the manifold is a state space on which we forecast where pathogens might go. MIRADOR adds a second, coupled manifold of immune receptors and treats it as both \emph{descriptive} (embedding natural repertoires) and \emph{prescriptive} (defining the feasible space of designed receptor panels). The mathematical ``mirror'' between these manifolds is the complementary co-embedding that allows us to talk about interception.

Operationally:

\begin{itemize}
  \item HERALD answers: \emph{Where is the antigen likely to drift?} (forecasting).
  \item MIRADOR answers: \emph{Where must a receptor panel live so that forecasted drift paths come within a binding radius?} (interception).
\end{itemize}

The remainder of this section makes these objects precise in a Davis-style language.

\subsection{Antigen manifold: Davis structure on viral drift}
\label{sec:antigen-manifold}

We begin by specializing the Davis construction to antigens.

\paragraph{Observation space and encoder.}
Let $\Acal$ denote the antigen space (e.g., viral spike sequences or sequence--structure pairs). An encoder
\[
  f_A : \Acal \to \mathbb{R}^d
\]
maps each antigen $a \in \Acal$ to an ambient representation $f_A(a) \in \mathbb{R}^d$. We define the antigen manifold as the image
\[
  \M_A := f_A(\Acal) \subset \mathbb{R}^d.
\]
On $\M_A$ we equip a Riemannian metric $g_A$ (for example, the pullback metric induced by $f_A$) with geodesic distance $d_{g_A}$. The ambient Euclidean distance is
\[
  \delta_A(z, z') := \bigl\| z - z' \bigr\|_2, \qquad z, z' \in \M_A.
\]
We assume $\M_A$ is endowed with Davis structure: along appropriate path families (below), geodesic and Euclidean distances remain within a bounded-distortion regime.

We write
\[
  u_A(a) := f_A(a) \in \M_A
\]
for the antigen state of $a \in \Acal$, and occasionally $z_A$ for a generic point in $\M_A$ when the underlying antigen is not explicit.

\paragraph{Configuration map and coarse labels.}
MIRADOR inherits the configuration/margin machinery of Davis manifolds. Concretely, we assume:

\begin{itemize}
  \item A configuration map
  \[
    c_A : \M_A \to \mathbb{R}^{p_A}
  \]
  summarizing neutralization-relevant features (e.g., cluster coordinates, antigenic principal components, epitope scores).

  \item A coarse label map
  \[
    h_A : \mathbb{R}^{p_A} \to \{\text{vax-prox},\text{escape},\text{amb}\},
  \]
  where $h_A \circ c_A$ partitions antigen states into vaccine-proximal, escape-like, and ambiguous regions.

  \item Hard and soft configuration margins $0 < \kappa^{(A)}_{\mathrm{hard}} \le \kappa^{(A)}_{\mathrm{soft}}$ that define an ambiguity band between vaccine-like and escape-like regions: inside $\kappa^{(A)}_{\mathrm{hard}}$ we are confident in the coarse label; between $\kappa^{(A)}_{\mathrm{hard}}$ and $\kappa^{(A)}_{\mathrm{soft}}$ the theory recommends abstention.
\end{itemize}

Let $R^{(A)}_0, R^{(A)}_1 \subset \M_A$ be non-empty closed reference sets summarizing canonical vaccine-proximal and escape-like configurations, respectively. We define the geodesic distance to each set by
\[
  d_{g_A}(z, R^{(A)}_y) := \inf_{z' \in R^{(A)}_y} d_{g_A}(z, z'), \qquad y \in \{0,1\}.
\]
The margins $(\kappa^{(A)}_{\mathrm{hard}}, \kappa^{(A)}_{\mathrm{soft}})$ are interpreted relative to these reference sets.

\begin{definition}[Antigen manifold]
\label{def:antigen-manifold}
A \emph{Davis antigen manifold} is a tuple $(\M_A, g_A, c_A, h_A, \kappa^{(A)}_{\mathrm{hard}}, \kappa^{(A)}_{\mathrm{soft}})$ consisting of an encoder $f_A : \Acal \to \mathbb{R}^d$ with image $\M_A = f_A(\Acal)$ equipped with pullback metric $g_A$, a configuration map $c_A : \M_A \to \mathbb{R}^{p_A}$, a coarse label map $h_A$, and margins $(\kappa^{(A)}_{\mathrm{hard}}, \kappa^{(A)}_{\mathrm{soft}})$ as described above, supplemented by a benign path family $\Pcal_A(L_A^\star)$ and distortion profile $\varepsilon_A(\cdot)$ satisfying the non-vacuity condition~\eqref{eq:antigen-nonvacuity}.
\end{definition}

\paragraph{Benign antigen paths.}
\begin{definition}[Benign antigen paths]
\label{def:benign-antigen-paths}
\label{def:mirror-paths}
A (piecewise) smooth path on $\M_A$ is a map $\Gamma : [0,1] \to \M_A$. Its geodesic length is
\[
  \ell_{g_A}(\Gamma) := \int_0^1 \bigl\|\dot{\Gamma}(t)\bigr\|_{g_A} \, dt.
\]
For $L > 0$, a family of \emph{benign antigen paths} $\Pcal_A(L)$ is a collection of such paths with
\[
  \ell_{g_A}(\Gamma) \le L \quad \text{for all } \Gamma \in \Pcal_A(L),
\]
where ``benign'' denotes biologically plausible, identity-preserving trajectories respecting structural viability and fitness constraints (e.g., epitope-restricted mutation paths).
\end{definition}

In practice, $\Pcal_A(L)$ might consist of epitope-restricted mutation paths with at most $L$ substitutions in a specified epitope set, or more general sequence/structure trajectories constrained by biophysical filters. The path horizon $L$ controls how far HERALD-style forecasting and MIRADOR-style interception are asked to operate.

\paragraph{Pathwise distortion profile and operational horizon.}
For a path $\Gamma \in \Pcal_A(L)$ and times $0 \le s < t \le 1$ with $d_{g_A}\bigl(\Gamma(s),\Gamma(t)\bigr) > 0$, define the distortion ratio
\[
  \rho^{(A)}_\Gamma(s,t)
  := \frac{\delta_A\bigl(\Gamma(s),\Gamma(t)\bigr)}{d_{g_A}\bigl(\Gamma(s),\Gamma(t)\bigr)}.
\]
The distortion radius at horizon $L$ is
\[
  \varepsilon_A(L)
  := \sup_{\Gamma \in \Pcal_A(L)}
     \;\sup_{\substack{0 \le s < t \le 1\\ d_{g_A}(\Gamma(s),\Gamma(t))>0}}
     \bigl| \rho^{(A)}_\Gamma(s,t) - 1 \bigr|.
\]
A bounded-distortion regime is one in which $\varepsilon_A(L)$ is small over the path horizon of interest. Distortion is audited empirically via such ratios on validation paths rather than computed exactly.

We assume the existence of an \emph{operational horizon} $L_A^\star > 0$ and tolerance $\varepsilon_0 > 0$ such that
\[
  \varepsilon_A(L_A^\star) < \varepsilon_0
\]
(e.g., $\varepsilon_0 = 0.1$), and distortions beyond $L_A^\star$ are treated as out-of-regime.

\paragraph{Non-vacuity of antigen margins.}
Let $R_A > 0$ characterize the geodesic scale of the antigen region of interest, for example any radius such that $R^{(A)}_0 \cup R^{(A)}_1$ is contained in a $d_{g_A}$-ball of radius $R_A$ and the system attempts to operate only within a $2R_A$ neighborhood of these references. A Davis-style non-vacuity condition for MIRADOR on the antigen side is
\begin{equation}
  \kappa^{(A)}_{\mathrm{soft}} - 2 R_A \,\varepsilon_A(L_A^\star) > 0,
  \label{eq:antigen-nonvacuity}
\end{equation}
ensuring that bounded geodesic--Euclidean distortion along benign paths does not completely erode the soft margin separating vaccine-proximal from escape-like regions.

\subsection{Receptor manifold: descriptive and prescriptive geometry}
\label{sec:receptor-manifold}

We now introduce an analogous Davis structure on receptors, with an important twist: the receptor manifold is both a description of natural repertoires and a design space for intervention.

\paragraph{Observation space and encoder.}
Let $\Rcal$ denote the receptor space (e.g., BCR or antibody sequences, possibly with structural annotations). An encoder
\[
  f_R : \Rcal \to \mathbb{R}^d
\]
maps each receptor $r \in \Rcal$ to an ambient representation $f_R(r) \in \mathbb{R}^d$. The receptor manifold is
\[
  \M_R := f_R(\Rcal) \subset \mathbb{R}^d,
\]
equipped with a Riemannian metric $g_R$ with geodesic distance $d_{g_R}$ and ambient distance
\[
  \delta_R(z, z') := \bigl\| z - z' \bigr\|_2, \qquad z, z' \in \M_R.
\]
As on the antigen side, Davis structure is assumed only on a path-restricted regime.

We write
\[
  u_R(r) := f_R(r) \in \M_R
\]
for the receptor state of $r \in \Rcal$.

\paragraph{Configuration map and coarse labels.}
The receptor manifold carries its own configuration structure, reflecting which regions correspond to viable, manufacturable, and immunologically reasonable receptors.

\begin{itemize}
  \item A configuration map
  \[
    c_R : \M_R \to \mathbb{R}^{p_R}
  \]
  summarizes receptor-relevant features (e.g., developability scores, paratope descriptors, isotype or Fc-engineering choices).

  \item A coarse label map
  \[
    h_R : \mathbb{R}^{p_R} \to \{\text{productive},\text{non-viable},\mathrm{amb}\},
  \]
  partitioning receptors into productive (developable, on-target), non-viable (e.g., unstable, polyreactive, or otherwise disallowed), and ambiguous regions.

  \item Hard and soft margins $0 < \kappa^{(R)}_{\mathrm{hard}} \le \kappa^{(R)}_{\mathrm{soft}}$ defining a receptor-side ambiguity band analogous to the antigen case.
\end{itemize}

Let $R^{(R)}_{\mathrm{prod}}, R^{(R)}_{\mathrm{bad}} \subset \M_R$ be non-empty closed reference sets summarizing canonical productive and non-viable receptor configurations, respectively, and define distances $d_{g_R}(z, R^{(R)}_{\cdot})$ analogously.

\begin{definition}[Receptor manifold]
\label{def:receptor-manifold}
A \emph{Davis receptor manifold} is a tuple $(\M_R, g_R, c_R, h_R, \kappa^{(R)}_{\mathrm{hard}}, \kappa^{(R)}_{\mathrm{soft}})$ consisting of an encoder $f_R : \Rcal \to \mathbb{R}^d$ with image $\M_R = f_R(\Rcal)$ equipped with pullback metric $g_R$, a configuration map $c_R : \M_R \to \mathbb{R}^{p_R}$, a coarse label map $h_R$, and margins $(\kappa^{(R)}_{\mathrm{hard}}, \kappa^{(R)}_{\mathrm{soft}})$ as described above, supplemented by a benign path family $\Pcal_R(L_R^\star)$ and distortion profile $\varepsilon_R(\cdot)$ satisfying the receptor-side non-vacuity condition.
\end{definition}

\paragraph{Benign receptor paths and design moves.}
\begin{definition}[Benign receptor paths]
\label{def:benign-receptor-paths}
A (piecewise) smooth receptor path is a map $\Gamma_R : [0,1] \to \M_R$ with geodesic length
\[
  \ell_{g_R}(\Gamma_R) := \int_0^1 \bigl\|\dot{\Gamma}_R(t)\bigr\|_{g_R} \, dt.
\]
For $L > 0$, a family of \emph{benign receptor paths} $\Pcal_R(L)$ consists of such paths with $\ell_{g_R}(\Gamma_R) \le L$ that approximate representative affinity-maturation trajectories (smooth approximations to the expected path under stochastic somatic hypermutation and selection) or computational design steps (e.g., rational substitutions, restricted library sampling, or gradient-guided moves in sequence space) subject to biological plausibility and feasibility constraints.
\end{definition}

As on the antigen side, we define a receptor-side distortion radius $\varepsilon_R(L)$ by taking suprema of distortion ratios $\rho^{(R)}_{\Gamma_R}(s,t)$ along $\Gamma_R \in \Pcal_R(L)$, and an operational horizon $L_R^\star$ with $\varepsilon_R(L_R^\star) < \varepsilon_0$ for some tolerance.

The path families $\Pcal_R(L)$ thus encode not only biological plausibility (e.g., maintaining structural stability, avoiding aggregation) but also practical design constraints (manufacturability, immunogenicity profiles, regulatory precedent), with horizon $L_R^\star$ defining the validated regime where such constraints remain active.

A receptor-side non-vacuity condition analogous to~\eqref{eq:antigen-nonvacuity} ensures that configuration margins on $\M_R$ are not destroyed by bounded distortion:
\[
  \kappa^{(R)}_{\mathrm{soft}} - 2 R_R\,\varepsilon_R(L_R^\star) > 0,
\]
for some radius $R_R > 0$ characterizing the receptor region of interest.

In contrast to HERALD, where the antigen manifold is purely \emph{observational}, the receptor manifold is both descriptive (embedding natural repertoires to understand existing responses) and prescriptive (defining the feasible space for designed receptors subject to manufacturability, immunogenicity, and safety constraints).

\subsection{Complementary co-embedding and lock-and-key semantics}
\label{sec:mirror-geometry}

MIRADOR couples $\M_A$ and $\M_R$ through a shared ambient space and an explicitly \emph{complementary} interpretation of distances between antigens and receptors.

\paragraph{Shared ambient space and shorthand notation.}
We assume $f_A$ and $f_R$ map into the same ambient $\mathbb{R}^d$, so that for $a \in \Acal$ and $r \in \Rcal$ we can define an antigen--receptor distance
\[
  \delta\bigl(a, r\bigr) := \bigl\|u_A(a) - u_R(r)\bigr\|_2.
\]
When a path $\Gamma$ on $\M_A$ is implicit, we will write
\[
  \delta\bigl(\Gamma(t), r\bigr)
\]
as shorthand for $\delta\bigl(u_A(\Gamma(t)), u_R(r)\bigr)$ to reduce clutter.

\paragraph{Complementary rather than similarity embeddings.}
The key MIRADOR assumption is that the shared embedding is \emph{lock-and-key} rather than similarity-based. Informally, $u_A(a)$ represents the antigen in antigen space, while $u_R(r)$ represents the location of the \emph{platonic antigen} that receptor $r$ would ideally bind.

Formally, for antigen--receptor pairs with measured or modeled binding data, we train $f_A$ and $f_R$ so that
\[
  \Phi_{\mathrm{bind}}\!\Bigl(\delta\bigl(a,r\bigr)\Bigr)
  \approx -\Delta G_{\mathrm{bind}}(a,r),
\]
where $\Delta G_{\mathrm{bind}}(a,r) = G_{\mathrm{complex}} - G_{\mathrm{free}}$ is the Gibbs free energy change upon binding (negative for favorable binding), and $\Phi_{\mathrm{bind}}$ is a decreasing calibration function so that smaller $\delta$ corresponds to more favorable (more negative) $\Delta G_{\mathrm{bind}}$.\footnote{In practice $\Phi_{\mathrm{bind}}$ can be learned from assay or simulation data and may be modeled as a simple monotone link (e.g., spline or isotonic regression) rather than a fixed analytic form.}
A binding-distance threshold is then defined as
\[
  \rho_{\mathrm{bind}} := \Phi_{\mathrm{bind}}^{-1}\bigl(\Delta G_{\mathrm{thresh}}\bigr),
\]
for some pre-registered free-energy cutoff $\Delta G_{\mathrm{thresh}}$ (e.g., $-10\,\mathrm{kcal/mol}$).

Crucially, the complementary structure is enforced via \emph{asymmetric contrastive training}: for each receptor $r$, we push $u_R(r)$ toward its known high-affinity partners in antigen space and away from three kinds of negatives:
\begin{enumerate}
  \item non-cognate antigens (no appreciable binding to $r$);
  \item antigens that bind other receptors but not $r$ (mismatched negatives);
  \item antigens that bind nothing in the panel (forcing specificity).
\end{enumerate}
This negative structure prevents trivial solutions in which all receptors collapse onto a single antigen centroid and is elaborated in the construction section (Section~\ref{sec:mirror-construction}).

\begin{definition}[Complementary co-embedding]
\label{def:complementary-coembedding}
Given a Davis antigen manifold $(\M_A,g_A)$ and receptor manifold $(\M_R,g_R)$ both embedded in $\mathbb{R}^d$, a \emph{complementary co-embedding} is the pair $(u_A,u_R)$ together with a decreasing calibration function $\Phi_{\mathrm{bind}}$ such that the antigen--receptor distance $\delta(a,r) := \|u_A(a)-u_R(r)\|_2$ satisfies $\Phi_{\mathrm{bind}}(\delta(a,r)) \approx -\Delta G_{\mathrm{bind}}(a,r)$, with binding radius $\rho_{\mathrm{bind}} := \Phi_{\mathrm{bind}}^{-1}(\Delta G_{\mathrm{thresh}})$ for a pre-registered free-energy cutoff $\Delta G_{\mathrm{thresh}}$. The embeddings are trained with asymmetric contrastive objectives so that $u_R(r)$ lies near the center of mass of its cognate antigen cluster.
\end{definition}

\paragraph{Mirror manifolds (terminology).}
The term ``mirror manifold'' in this work is purely metaphorical: it denotes the complementary, co-embedded pair $(\M_A,\M_R)$ together with the distance-based lock-and-key semantics above. It is not related to mirror symmetry in string theory or Calabi--Yau duals.

\subsection{Interception as path-wise coverage}
\label{sec:interception-coverage}

With the antigen and receptor manifolds coupled via a complementary co-embedding, MIRADOR defines interception geometrically as path-wise coverage.

\paragraph{Forecasted drift and path distribution.}
Let $\Pi_A(L)$ denote a forecast distribution over benign antigen drift paths $\Gamma : [0,1] \to \M_A$ drawn from $\Pcal_A(L)$ with $L \le L_A^\star$. This distribution encodes HERALD-style beliefs about where a lineage may drift over a given horizon (e.g., a season or other time window), restricted to the Davis-valid regime.

\begin{definition}[Hit event and coverage]
\label{def:hit-event}
Given a candidate receptor panel $V \subset \Rcal$ and antigen path $\Gamma : [0,1] \to \M_A$ drawn from $\Pi_A(L)$, the \emph{hit} event is
\[
  \Hit(V,\Gamma)
  := \Bigl\{\exists\,t \in [0,1],\ \exists\,r \in V :
         \delta\bigl(\Gamma(t), r\bigr) \le \rho_{\mathrm{bind}} \Bigr\},
\]
i.e., the path passes within the binding radius of at least one receptor in $V$ at some time along its evolution. The \emph{coverage} of $V$ under the forecast distribution is
\[
  C(V) := \Pr_{\Gamma \sim \Pi_A(L)}\bigl(\Hit(V,\Gamma)\bigr) \in [0,1].
\]
In later sections we will distinguish $C(V)$ from empirical estimates $\widehat{C}(V)$ obtained by sampling a finite number of drift paths from $\Pi_A(L)$ and discretizing them into representative antigenic states.
\end{definition}

\paragraph{Geometric set cover viewpoint.}
At a computational level, the path-wise definition above induces a finite set cover problem on a Riemannian manifold. Sampling $\Gamma^{(1)},\dots,\Gamma^{(N)} \sim \Pi_A(L)$ and discretizing each path into antigenic states $\{a_i\}_{i=1}^N$ with sufficient density to ensure $|\widehat{C}(V) - C(V)| \le \varepsilon_{\mathrm{samp}}$ with high probability, the design problem becomes:

\begin{quote}
  Choose a panel $V$ of receptors such that each sampled antigen $a_i$ lies within $\rho_{\mathrm{bind}}$ of at least one $r \in V$, while minimizing panel size or maximizing coverage subject to a size budget.
\end{quote}

This is a geometric set cover problem on $(\M_A,g_A)$: the ``sets'' are receptor-centered balls of radius $\rho_{\mathrm{bind}}$ under the complementary metric, and the ``elements'' are samples from forecast drift paths. MIRADOR leverages standard approximation algorithms for set cover on finite sets (e.g., greedy selection) and then analyzes how their performance lifts back to the underlying continuous drift family once geometric distortion, sampling error, and model uncertainty are folded into an inverted error budget in Section~\ref{sec:inverted-error-budgets}.

\paragraph{Design goals and coverage certificates.}
In summary, MIRADOR's design goal is to construct receptor panels $V$ that:
\begin{enumerate}
  \item \emph{achieve high coverage} $C(V)$ against forecasted drift families of interest; and
  \item \emph{provide a coverage certificate}: a lower bound on $C(V)$ that accounts for geometric distortion, binding-model uncertainty, sampling error, and linkage error between binding and neutralization.
\end{enumerate}
The complementary manifolds $(\M_A,\M_R)$ and the interception event $\Hit(V,\Gamma)$ defined in this section provide the geometric substrate on which those certificates will be constructed in subsequent sections.


\section{Mirror Construction: Training the Receptor Manifold and Co-Embedding}
\label{sec:mirror-construction}

Section~\ref{sec:mirror-geometry} defined the antigen and receptor manifolds
$(\M_A,g_A)$, $(\M_R,g_R)$, the complementary co-embedding
$(u_A,u_R)$ of Definition~\ref{def:complementary-coembedding}, and mirror
path families of Definition~\ref{def:mirror-paths}. In this section we give a
concrete construction of the receptor side of the mirror geometry:
we specify training data, loss functions, negative sampling, smoothness
regularization, and path-generation procedures that jointly produce

\begin{enumerate}
  \item a Davis-style receptor manifold $(\M_R,g_R,\Pcal_R(L_R^\star),\varepsilon_R(\cdot))$ with bounded distortion along benign receptor paths; and
  \item a complementary co-embedding $(u_A,u_R)$ whose receptor coordinates
  behave as “lock–and–key” centers of mass for cognate antigens on a binding
  in-distribution region.
\end{enumerate}

We work throughout with the antigen manifold $(\M_A,g_A)$ already constructed
in Section~\ref{sec:mirror-geometry}, and treat the antigen embedding
$u_A : \Acal \to \M_A$ as fixed.

\subsection{Binding data, cognate sets, and negative pools}

Let $\Acal$ denote the antigen space and $\Rcal$ the receptor design space
(e.g., antibody or receptor sequences that satisfy basic manufacturability and
safety constraints). Let
\[
  \Dcal_{\mathrm{bind}} \subseteq \Acal \times \Rcal \times \{0,1\}
\]
be a binding dataset of labeled pairs $(a,r,y)$, where $y=1$ denotes
“binds above a pre-registered affinity threshold” and $y=0$ denotes
“does not bind above that threshold.”
We write $(a,r)\in\Dcal_{\mathrm{bind}}^+$ when $y=1$ and
$(a,r)\in\Dcal_{\mathrm{bind}}^-$ when $y=0$.

For each receptor $r \in \Rcal$ we define its observed cognate and non-cognate
antigen sets
\begin{align*}
  \Acal_r^+ &:= \{ a \in \Acal : (a,r,1) \in \Dcal_{\mathrm{bind}} \},\\
  \Acal_r^{\mathrm{nc}} &:= \{ a \in \Acal : (a,r,0) \in \Dcal_{\mathrm{bind}} \}.
\end{align*}

We assume that a subset of receptors has at least one observed positive:

\begin{equation}
  \Rcal_{\mathrm{train}} := \bigl\{ r \in \Rcal : |\Acal_r^+| \ge 1 \bigr\},
  \qquad
  \Acal_r^+ \neq \emptyset \text{ for all } r \in \Rcal_{\mathrm{train}}.
  \label{eq:Rtrain}
\end{equation}

Receptors outside $\Rcal_{\mathrm{train}}$ may still be embedded by $u_R$ using
the base and smoothness losses only (Section~\ref{subsec:receptor-loss}),
but they are excluded from the center-of-mass and contrastive terms and are
not used for coverage certificates.\footnote{Receptors with
$\Acal_r^+ = \emptyset$ can still be embedded by $u_R$ (via $L_{\mathrm{base}}$
and $L_{\mathrm{smooth}}^{(R)}$), for example to ensure geometric regularity or
to support other tasks, but they do not contribute to the lock–and–key
objective or the InfoNCE loss, and panels used for MIRADOR coverage
certificates are restricted to $V \subseteq \Ical^R$ as defined below.}

We further define two global negative pools in antigen space:
\begin{align*}
  \Acal_r^{\mathrm{mm}}
    &:= \Bigl( \bigcup_{r' \in \Rcal_{\mathrm{train}}\setminus\{r\}} \Acal_{r'}^+ \Bigr)
        \setminus \Acal_r^+ 
        &&\text{(mismatched positives: bind other receptors, not $r$)},\\
  \Acal^{\mathrm{null}}
    &:= \Bigl\{ a \in \Acal : \forall r' \in \Rcal_{\mathrm{train}},\ (a,r',1)\notin\Dcal_{\mathrm{bind}} \Bigr\}
        &&\text{(null antigens: bind no training receptor)}.
\end{align*}

These three sets play distinct roles in the contrastive loss:

\begin{itemize}
  \item $\Acal_r^{\mathrm{nc}}$ (non-cognate for $r$): “same receptor, non-binding antigen.”
  \item $\Acal_r^{\mathrm{mm}}$ (mismatched): “antigens that bind something else, not $r$.”
  \item $\Acal^{\mathrm{null}}$ (null): “antigens that bind nothing in the panel.”
\end{itemize}

Given a receptor $r\in\Rcal_{\mathrm{train}}$ we define a negative-sampling
mixture
\begin{equation}
  Q_{\mathrm{neg}}(\cdot \mid r)
    := w_{\mathrm{nc}}\, U(\Acal_r^{\mathrm{nc}})
     + w_{\mathrm{mm}}\, U(\Acal_r^{\mathrm{mm}})
     + w_{\mathrm{null}}\, U(\Acal^{\mathrm{null}}),
  \label{eq:Qneg}
\end{equation}
where $U(\cdot)$ denotes the uniform distribution over its argument and
$w_{\mathrm{nc}},w_{\mathrm{mm}},w_{\mathrm{null}}\ge 0$ with
$w_{\mathrm{nc}}+w_{\mathrm{mm}}+w_{\mathrm{null}}=1$. In practice we draw a
mini-batch of negatives $\{a_k^-\}_{k=1}^K$ (e.g.\ $K \in [10,100]$) from
$Q_{\mathrm{neg}}(\cdot\mid r)$, either uniformly over the three pools or with
user-chosen weights $(w_{\mathrm{nc}},w_{\mathrm{mm}},w_{\mathrm{null}})$ to emphasize
particular negative types.

\subsection{Binding in-distribution region}
\label{subsec:binding-indist}

Let $\delta_A(a,a') := \|u_A(a)-u_A(a')\|_2$ denote the Euclidean distance in
antigen space (Section~\ref{sec:mirror-geometry}), which approximates geodesic
distance $d_{g_A}$ in the bounded-distortion regime.

\begin{definition}[Binding in-distribution region]
\label{def:binding-in-distribution}
Fix a coverage radius $\eta_A>0$ in antigen space. For each receptor
$r \in \Rcal_{\mathrm{train}}$, define the antigenic in-distribution neighborhood
\[
  \Ical_r^A
    := \Bigl\{ a \in \Acal :
         \exists (a',r') \in \Dcal_{\mathrm{bind}}
         \ \text{with}\ 
         \delta_A(a,a') \le \eta_A \Bigr\}.
\]
The global antigenic in-distribution region for binding is
\[
  \Ical^A := \bigcup_{r \in \Rcal_{\mathrm{train}}} \Ical_r^A.
\]
We take the receptor in-distribution region for binding to be
\[
  \Ical^R := \Rcal_{\mathrm{train}},
\]
so that only receptors with at least one observed cognate antigen are used for
lock–and–key training and coverage certificates. The joint binding
in-distribution region is
\[
  \Omega_{\mathrm{bind}} := \Ical^A \times \Ical^R.
\]
\end{definition}

Informally, $\Ical^A$ collects antigenic states lying within a small radius
$\eta_A$ of antigens observed in binding assays, and $\Ical^R$ collects
receptors with sufficient coverage. All guarantees in this section are intended
to hold on $\Omega_{\mathrm{bind}}$; receptors or antigens far from
$\Omega_{\mathrm{bind}}$ are treated as out-of-distribution for binding and are
handled by abstention or conservative heuristics.

\subsection{Lock--key interpretation and center-of-mass loss}

Recall from Section~\ref{sec:mirror-geometry} that $u_A(a) := f_A(a)\in\M_A$
and $u_R(r) := f_R(r)\in\M_R$ denote the antigen and receptor embeddings,
abusing notation by identifying these with their images in $\R^d$ via the
coordinate maps. We write $\delta(a,r)$ as shorthand for
$\delta(u_A(a),u_R(r))$.

\paragraph{Cognate distributions.}
For each receptor $r\in\Rcal$ we posit an underlying (typically unknown)
cognate antigen distribution $A_r$ over $\Acal$, defined as the distribution of
antigens that bind $r$ above a pre-registered affinity threshold
$\Delta G_{\mathrm{thresh}}$:
\[
  A_r(a) \propto \Pr\bigl( \Delta G(a,r) \le \Delta G_{\mathrm{thresh}} \bigr)
  \quad\text{for } a \in \Acal,
\]
where $\Delta G(a,r)$ is the physical binding free energy. In practice we only
observe a finite set $\Acal_r^+$ of positives and approximate $A_r$ by its
empirical version supported on $\Acal_r^+$.

Let $w_r(a) \ge 0$ for $a\in\Acal_r^+$ be optional importance weights
(e.g.\ reflecting assay quality or diversity), normalized so that
$\sum_{a\in\Acal_r^+} w_r(a)=1$. The empirical center-of-mass of cognate
antigens in embedding space is
\begin{equation}
  \hat{\mu}_r
    := \sum_{a\in\Acal_r^+} w_r(a)\, u_A(a)
    \in \M_A.
  \label{eq:mu-hat}
\end{equation}

\begin{definition}[Lock--key interpretation]
For each receptor $r\in\Rcal_{\mathrm{train}}$, the mirror geometry is intended
to satisfy
\[
  u_R(r) \approx \E_{a \sim A_r}\bigl[u_A(a)\bigr],
\]
i.e.\ $u_R(r)$ approximates the center-of-mass of its cognate antigen
distribution in embedding space. In practice this expectation is approximated
by $\hat{\mu}_r$ from \eqref{eq:mu-hat} using observed cognate antigens
$\Acal_r^+$.
\end{definition}

This interpretation motivates a center-of-mass loss
\begin{equation}
  L_{\mathrm{cm}}(r;\theta_R)
    := \bigl\| u_R(r) - \hat{\mu}_r \bigr\|_2^2,
  \qquad r \in \Rcal_{\mathrm{train}},
  \label{eq:Lcm}
\end{equation}
which encourages $u_R(r)$ to sit near the average of its cognate antigen
embeddings.

\subsection{Asymmetric contrastive training with negative mixtures}
\label{subsec:contrastive-R}

We couple the center-of-mass objective with an asymmetric contrastive loss that
uses the three negative pools in~\eqref{eq:Qneg} to prevent trivial collapse
solutions.

For a receptor--antigen pair $(r,a)$ we define a similarity score
\begin{equation}
  s_{\theta_R}(r,a)
    := -\frac{\bigl\|u_R(r) - u_A(a)\bigr\|_2^2}{2\,\tau_{\mathrm{bind}}^2},
  \label{eq:binding-score}
\end{equation}
where $\tau_{\mathrm{bind}}>0$ is a temperature parameter that controls the
scale of the contrastive gradients. In practice $\tau_{\mathrm{bind}}$ can be
tuned via cross-validation or set to a default (e.g.\ $\tau_{\mathrm{bind}}=0.1$)
that balances gradient magnitude and numerical stability.

For each positive pair $(r,a^+)$ with $a^+ \in \Acal_r^+$ we draw a mini-batch
of negatives $\{a_k^-\}_{k=1}^K$ from $Q_{\mathrm{neg}}(\cdot\mid r)$ and
define a receptor–antigen InfoNCE loss
\begin{equation}
  L_{\mathrm{NCE}}^{(R)}(\theta_R)
  := \E_{(r,a^+) \sim \Dcal_{\mathrm{bind}}^+}
     \Biggl[
       -\log
       \frac{\exp\bigl\{s_{\theta_R}(r,a^+)\bigr\}}
            {\exp\bigl\{s_{\theta_R}(r,a^+)\bigr\}
              + \sum_{k=1}^K
                   \exp\bigl\{s_{\theta_R}(r,a_k^-)\bigr\}}
     \Biggr].
  \label{eq:InfoNCE-R}
\end{equation}

As in standard contrastive learning, small values of
$L_{\mathrm{NCE}}^{(R)}(\theta_R)$ imply that, on the in-distribution region,
positive (cognate) pairs receive higher similarity scores and smaller distances
than negatives by a positive margin, which induces a distance gap in the
co-embedding.

To make this explicit we define the expected binding distance gap on the
(binding) in-distribution region $\Omega_{\mathrm{bind}}$:

\begin{equation}
  \Delta_{\mathrm{bind}}
    := \inf_{r \in \Ical^R}
      \Bigl(
        \E_{a \in \Acal_r^+} \bigl[ \delta(a,r) \bigr]
        - \E_{a \sim Q_{\mathrm{neg}}(\cdot\mid r)}
           \bigl[ \delta(a,r) \bigr]
      \Bigr).
  \label{eq:Delta-bind}
\end{equation}

We will assume in Theorem~\ref{thm:receptor-existence} that the trained
co-embedding achieves a strictly positive gap $\Delta_{\mathrm{bind}}>0$ on
$\Omega_{\mathrm{bind}}$.

\subsection{Receptor smoothness and benign receptor paths}
\label{subsec:receptor-paths}

We equip the receptor embedding with a smoothness regularizer
$L_{\mathrm{smooth}}^{(R)}(\theta_R)$ analogous to the Davis manifold
construction. On the receptor input space (e.g.\ sequence space) we consider
either a Jacobian energy
\begin{equation}
  L_{\mathrm{smooth}}^{(R)}(\theta_R)
    := \E_{r \sim \mu_R}
          \bigl\| \nabla_r u_R(r) \bigr\|_F^2
  \quad\text{for some reference measure $\mu_R$},
  \label{eq:smooth-jacobian}
\end{equation}
or a graph-Laplacian smoother
\begin{equation}
  L_{\mathrm{smooth}}^{(R)}(\theta_R)
    := \mathrm{Tr}\bigl( F_R^\top L_R F_R \bigr),
  \quad
  F_R = [u_R(r_1),\dots,u_R(r_n)]^\top,
  \label{eq:smooth-laplacian}
\end{equation}
where $L_R$ is a $k$NN Laplacian over receptors (or receptor sequences)
and $\{r_i\}$ are sampled receptors.

As in the Davis existence theorem, increasing the smoothness weight
$\lambda_R$ shrinks curvature of the pullback metric $g_R$ induced by $u_R$
on $\M_R$ and yields a receptor distortion profile
$\varepsilon_R(L)$ along benign receptor paths
$\Pcal_R(L)$ satisfying
\[
  \varepsilon_R(L) \lesssim C_R' K_R(\lambda_R) L,
\]
for $L \le L_R^\star$ and some curvature proxy $K_R(\lambda_R)$ that decreases
with $\lambda_R$.

\paragraph{Generating receptor paths for auditing.}
In practice, receptor paths for distortion audits and for defining benign
families $\Pcal_R(L)$ can be constructed from three complementary sources:

\begin{enumerate}
  \item \emph{Lineage-based paths.} When affinity-maturation lineages are
    available (germline $\to$ intermediate $\to$ mature), we connect embedded
    states $u_R(r^{(0)}), u_R(r^{(1)}),\dots$ by geodesic interpolants in
    $\M_R$ and treat the resulting piecewise-geodesic curves as elements of
    $\Pcal_R(L)$.
  \item \emph{Computational design paths.} For sequence-level design, we form
    paths by chaining small-step mutations that preserve predicted stability
    and developability scores (e.g.\ no aggregation flags, acceptable
    biophysical heuristics), and we embed these as curves in $\M_R$.
  \item \emph{Random-walk paths.} We sample short Riemannian random walks
    (e.g.\ Brownian bridges) on $\M_R$ starting and ending at productive
    receptors, filtering out points that leave the productive configuration
    region under the receptor configuration map $h_R$.
\end{enumerate}

Distortion along these paths is audited empirically by comparing geodesic and
Euclidean lengths, as in the Davis framework; we retain only paths for which
the empirical receptor distortion radius $\varepsilon_R(L_R^\star)$ remains
below a pre-registered tolerance.

\subsection{Composite receptor loss}
\label{subsec:receptor-loss}

Putting the pieces together, we train the receptor encoder with a composite
loss
\begin{equation}
  L_R(\theta_R)
    := \alpha_R L_{\mathrm{base}}(\theta_R)
     + \beta_R L_{\mathrm{NCE}}^{(R)}(\theta_R)
     + \gamma_R \E_{r \in \Rcal_{\mathrm{train}}}
                 \bigl[ L_{\mathrm{cm}}(r;\theta_R) \bigr]
     + \lambda_R L_{\mathrm{smooth}}^{(R)}(\theta_R),
  \label{eq:receptor-loss}
\end{equation}
where:

\begin{itemize}
  \item $L_{\mathrm{base}}$ is an \emph{optional} task-specific loss (which may
    be zero if no auxiliary supervision is available), e.g.\ predicting paratope
    class, developability scores, or coarse receptor configuration labels;
  \item $L_{\mathrm{NCE}}^{(R)}$ is the receptor–antigen InfoNCE loss
    in~\eqref{eq:InfoNCE-R};
  \item $L_{\mathrm{cm}}$ is the lock–and–key center-of-mass loss
    in~\eqref{eq:Lcm};
  \item $L_{\mathrm{smooth}}^{(R)}$ is a smoothness penalty of the form
    \eqref{eq:smooth-jacobian} or \eqref{eq:smooth-laplacian}.
\end{itemize}

Weights $(\alpha_R,\beta_R,\gamma_R,\lambda_R)$ control the trade-offs between
task supervision, binding separation, lock–and–key structure, and geometric
regularity. We are particularly interested in regimes where $\beta_R$ and
$\gamma_R$ are large enough to enforce a positive binding margin
$\Delta_{\mathrm{bind}}>0$ on $\Omega_{\mathrm{bind}}$, and $\lambda_R$ is large
enough to keep $\varepsilon_R(L_R^\star)$ small.

\subsection{Center-of-mass behavior and collapse avoidance (informal)}
\label{subsec:cm-informal}

The next proposition summarizes the intended qualitative behavior of the
composite loss~\eqref{eq:receptor-loss}. It is stated informally, as a
guiding principle rather than a formal convergence theorem.

\begin{proposition}[Center-of-mass and collapse avoidance (informal)]
\label{prop:center-of-mass}
Assume:

\begin{enumerate}
  \item[(i)] For each receptor $r \in \Rcal_{\mathrm{train}}$, the cognate set
    $\Acal_r^+$ is drawn from a distribution $A_r$ with finite second moment
    and non-degenerate support in $\M_A$.
  \item[(ii)] The antigen embedding $u_A$ is approximately isometric on
    $\Ical^A$ in the sense of Section~\ref{sec:mirror-geometry}, so that
    Euclidean averages of $u_A(a)$ are meaningful summaries of $A_r$.
  \item[(iii)] The negative pools in~\eqref{eq:Qneg} are sufficiently separated
    from $\Acal_r^+$ so that, in the limit of infinite data, standard InfoNCE
    analysis yields a positive binding margin $\Delta_{\mathrm{bind}}>0$ on
    $\Omega_{\mathrm{bind}}$.
  \item[(iv)] The smoothness weight $\lambda_R$ is chosen so that the receptor
    distortion profile $\varepsilon_R(L)$ remains below a pre-registered
    tolerance on $\Pcal_R(L_R^\star)$.
\end{enumerate}

Then, under regularity conditions (bounded curvature of $(\M_R,g_R)$, sufficient
capacity, and suitable optimization), any well-optimized solution
$\theta_R^\star$ to~\eqref{eq:receptor-loss} approximately satisfies:

\begin{enumerate}
  \item $u_R(r) \approx \hat{\mu}_r$ for each $r \in \Rcal_{\mathrm{train}}$,
    i.e.\ the lock–and–key center-of-mass constraint is enforced;
  \item $u_R$ does not collapse receptors to a single antigen embedding:
    receptors with disjoint cognate sets $\Acal_r^+$ map to distinct regions of
    $\M_R$, separated by a distance at least proportional to the binding margin
    $\Delta_{\mathrm{bind}}$.
\end{enumerate}

In particular, the combination of the center-of-mass loss and the asymmetric
negative structure in $L_{\mathrm{NCE}}^{(R)}$ prevents trivial solutions where
all receptors map to a single antigen centroid.
\end{proposition}

A more formal statement could be derived by adapting standard InfoNCE
margin-to-separation arguments and convergence analyses for center-of-mass
objectives under quadratic losses; we leave this to future work and use
Proposition~\ref{prop:center-of-mass} as an informal guide.

\subsection{Existence of a receptor mirror manifold on the binding in-distribution region}
\label{subsec:existence-R}

We now state a receptor-side existence theorem that mirrors the Davis
existence result: under suitable conditions, minimizing the composite loss
$L_R$ yields a receptor manifold and complementary co-embedding that satisfy
the mirror-geometry conditions of Definition~\ref{def:complementary-coembedding}
and the mirror-path conditions of Definition~\ref{def:mirror-paths} on the
(binding) in-distribution region.

\begin{theorem}[Existence of receptor mirror geometry under composite training]
\label{thm:receptor-existence}
Fix antigen and receptor path horizons $L_A^\star>0$, $L_R^\star>0$ and radii
$R_A,R_R>0$ describing the regions of interest around the antigen and receptor
reference sets (as in Section~\ref{sec:mirror-geometry}). Assume:

\begin{enumerate}
  \item[\textnormal{(1)}] \textbf{Antigen-side Davis structure.}
    The antigen manifold $(\M_A,g_A)$, path family $\Pcal_A(L_A^\star)$, and
    distortion profile $\varepsilon_A(L)$ satisfy the Davis manifold conditions
    of Definition~\ref{def:antigen-manifold} on a radius-$R_A$ region containing
    all relevant antigen states, with
    $\varepsilon_A(L_A^\star) < \varepsilon_A^{\max}$ and a non-vacuous antigen
    margin in the sense that
    \[
      \kappa_{\mathrm{soft}}^{(A)} - 2R_A \varepsilon_A(L_A^\star) > 0.
    \]

  \item[\textnormal{(2)}] \textbf{Receptor smoothness and distortion.}
    The receptor embedding $u_R$ is trained with $L_{\mathrm{smooth}}^{(R)}$
    such that the receptor distortion profile satisfies
    \[
      \varepsilon_R(L) \le C_R'' K_R(\lambda_R) L
      \quad\text{for } L \le L_R^\star
    \]
    along benign receptor paths $\Pcal_R(L)$ in a radius-$R_R$ region
    containing $\Ical^R$, with $K_R(\lambda_R)$ decreasing in $\lambda_R$.
    Moreover, the receptor non-vacuity condition holds:
    \[
      \kappa_{\mathrm{soft}}^{(R)} - 2 R_R \varepsilon_R(L_R^\star) > 0
    \]
    for some receptor-side configuration margins
    $(\kappa_{\mathrm{hard}}^{(R)},\kappa_{\mathrm{soft}}^{(R)})$ defined as in
    Section~\ref{sec:mirror-geometry}.

  \item[\textnormal{(3)}] \textbf{Binding margin on the in-distribution region.}
    The composite loss~\eqref{eq:receptor-loss} is optimized so that, on the
    binding in-distribution region $\Omega_{\mathrm{bind}}$, the binding
    distance gap~\eqref{eq:Delta-bind} satisfies
    \[
      \Delta_{\mathrm{bind}} > 0,
    \]
    i.e.\ for all $r \in \Ical^R$,
    \begin{equation}
      \E_{a \in \Acal_r^+}\bigl[ \delta(a,r) \bigr]
        + \Delta_{\mathrm{bind}}
      \;\le\;
      \E_{a \sim Q_{\mathrm{neg}}(\cdot\mid r)}
        \bigl[ \delta(a,r) \bigr].
      \label{eq:positive-margin}
    \end{equation}
    In addition, the center-of-mass term is small relative to the intrinsic
    variance of cognate antigen embeddings, in the sense that for all
    $r \in \Ical^R$
    \[
      \bigl\| u_R(r) - \hat{\mu}_r \bigr\|_2^2
        \ll \Var_{a \sim A_r}\bigl[ u_A(a) \bigr].
    \]

  \item[\textnormal{(4)}] \textbf{Path support and in-distribution coverage.}
    Antigen drift and benign receptor paths are supported in the binding
    in-distribution region in the following sense:
    \begin{enumerate}
      \item[(a)] For $\Gamma_A \sim \Pi_A(L_A^\star)$, we have
        $\Gamma_A(t) \in u_A(\Ical^A)$ for all $t \in [0,1]$ with probability
        at least $1 - \varepsilon_{\mathrm{path}}$, for some small
        $\varepsilon_{\mathrm{path}} \in [0,1)$; and
      \item[(b)] For every $\Gamma_R \in \Pcal_R(L_R^\star)$ we have
        $\Gamma_R(t) \in u_R(\Ical^R)$ for all $t \in [0,1]$.
    \end{enumerate}
\end{enumerate}

Then, for sufficiently large smoothness weight $\lambda_R$ and sufficiently
large contrastive and center-of-mass weights $(\beta_R,\gamma_R)$, there exists
a trained receptor encoder $u_R$ such that:

\begin{enumerate}
  \item[\textnormal{(i)}] The receptor manifold $(\M_R,g_R)$ together with
    $\Pcal_R(L_R^\star)$ and $\varepsilon_R(\cdot)$ satisfies the Davis
    manifold conditions of Definition~\ref{def:receptor-manifold} on the
    radius-$R_R$ region containing $\Ical^R$, with non-vacuous receptor
    margins $(\kappa_{\mathrm{hard}}^{(R)},\kappa_{\mathrm{soft}}^{(R)})$.

  \item[\textnormal{(ii)}] On the binding in-distribution region
    $\Omega_{\mathrm{bind}} = \Ical^A \times \Ical^R$, the co-embedding
    $(u_A,u_R)$ satisfies the complementary co-embedding conditions of
    Definition~\ref{def:complementary-coembedding} with lock–and–key semantics:
    $u_R(r)$ lies near the center-of-mass of cognate antigens for each
    $r \in \Ical^R$, and the binding margin~\eqref{eq:positive-margin}
    separates cognate from non-cognate antigens.

  \item[\textnormal{(iii)}] For any panel $V \subseteq \Ical^R$ the hit event
    $\Hit(V,\Gamma_A)$ and coverage $C(V;L_A^\star)$ defined in
    Section~\ref{sec:mirror-geometry} are well-posed on the support of
    $\Pi_A(L_A^\star)$.

  \item[\textnormal{(iv)}] These geometric and binding quantities can be
    computed using Euclidean surrogates and the calibrated binding threshold
    $\rho_{\mathrm{bind}} = \Phi_{\mathrm{bind}}^{-1}(\Delta G_{\mathrm{thresh}})$
    in a regime where: (a) geodesic distances are well approximated by ambient
    Euclidean distances along benign paths (controlled by
    $\varepsilon_A(L_A^\star)$ and $\varepsilon_R(L_R^\star)$); and (b) the
    binding model is empirically supported by data on $\Omega_{\mathrm{bind}}$.
\end{enumerate}
\end{theorem}

\begin{proof}[Proof sketch]
Assumption~(1) provides an antigen-side Davis manifold with a non-vacuous
margin, as in the general Davis existence theorem. Assumption~(2), together
with the smoothness regularizer, yields a receptor distortion profile
$\varepsilon_R(L)$ that is linear in $L$ with slope controlled by
$K_R(\lambda_R)$, and allows $\lambda_R$ to be chosen so that the receptor
non-vacuity condition holds at horizon $L_R^\star$.

Assumption~(3) combines the InfoNCE term $L_{\mathrm{NCE}}^{(R)}$ and the
center-of-mass term $L_{\mathrm{cm}}$ to enforce both a positive binding margin
$\Delta_{\mathrm{bind}}$ and lock–and–key alignment on $\Omega_{\mathrm{bind}}$,
mirroring Proposition~\ref{prop:center-of-mass}. This yields receptor-side
configuration margins $(\kappa_{\mathrm{hard}}^{(R)},\kappa_{\mathrm{soft}}^{(R)})$
and prevents collapse to a single antigen centroid.

Assumption~(4) guarantees that the antigen drift paths of interest and benign
receptor paths remain in the binding in-distribution region, so that both the
antigen and receptor distortion bounds apply and the binding model is
supported by data. Combining the antigen and receptor Davis structures with
the binding margin and center-of-mass behavior verifies the items in the
definition of a complementary co-embedding on $\Omega_{\mathrm{bind}}$, and
ensures that hit events and coverage computations can be carried out using
Euclidean surrogates within the regime of bounded distortion. A full proof
would adapt the arguments of the Davis manifold existence theorem to the joint
antigen–receptor setting; we omit the technical details here.
\end{proof}

\paragraph{Discussion.}
Theorem~\ref{thm:receptor-existence} ties together the pieces of MIRADOR's
mirror construction: InfoNCE with asymmetric negatives, a lock–and–key
center-of-mass objective, and Davis-style smoothness regularization on the
receptor side. Under the stated assumptions, these ingredients produce a
receptor manifold and complementary co-embedding that:

\begin{itemize}
  \item preserve antigen-side Davis structure on the binding in-distribution
    region;
  \item endow receptor coordinates with a meaningful lock–and–key semantics;
  \item maintain bounded geodesic–Euclidean distortion along both antigen and
    receptor benign paths; and
  \item support well-posed hit and coverage computations for panels
    $V \subseteq \Ical^R$ in a regime where Euclidean computations are
    geometrically interpretable and supported by data.
\end{itemize}

The next section leverages this construction to build coverage certificates and
compositional error budgets for MIRADOR panels.

\section{Applications: Vaccine and Antibody Design on MIRADOR}
\label{sec:mirador-applications}

We now illustrate how MIRADOR’s mirror geometry and coverage certificates can be used to
formulate concrete design problems for vaccines and therapeutic antibodies. Throughout, we work
with the antigen manifold $(\mathcal{M}_A,g_A)$, the receptor manifold $(\mathcal{M}_R,g_R)$, the
complementary co-embedding $u_A,u_R$, the antigen path family
$\mathcal{P}_A(L_A^\star)$ at operational horizon $L_A^\star$, and the coverage certificates
$\underline{C}_{\mathrm{neut}}(V;L_A^\star)$ derived from the inverted error-budget framework.

\subsection{Geometric formulation of vaccine design}
\label{subsec:geo-vaccine-design}

\paragraph{Vaccine footprints as receptor-like anchors.}
Consider a candidate vaccine payload or immunogen $\nu$ (e.g., an mRNA-encoded spike or a
cocktail of antigens). Conceptually, $\nu$ induces a polyclonal antibody repertoire
$\mathcal{R}_\nu \subseteq \mathcal{R}$, which MIRADOR represents in mirror space via
$u_R(\mathcal{R}_\nu) \subseteq \mathcal{M}_R$. For design, we summarize this repertoire by a
finite \emph{vaccine footprint}
\[
  V_\nu := \{v_1,\dots,v_K\} \subseteq \mathcal{I}^R \subseteq \mathcal{M}_R,
\]
where $\mathcal{I}^R$ is the (binding) in-distribution region on the receptor side, and each
$v_k$ is a receptor-like anchor that approximates a cluster of elicited neutralizing antibodies.
A multivalent vaccine or schedule corresponds to a union of such footprints.

\paragraph{Strong path-wise geometric coverage.}
Let $\Gamma_A \in \mathcal{P}_A(L_A^\star)$ be an antigenic path, and let
$\delta(a,r) := \delta\bigl(u_A(a),u_R(r)\bigr)$ denote the ambient mirror-space distance.
For a finite footprint $V \subseteq \mathcal{I}^R$ define the worst-case mirror distance along
$\Gamma_A$,
\begin{equation}
  d_{\max}(\Gamma_A,V)
  := \sup_{t \in [0,1]} \;\min_{v \in V} \delta\bigl(\Gamma_A(t),v\bigr).
\end{equation}

\begin{definition}[Strong mirror-space coverage of a path]
\label{def:strong-path-coverage}
Fix a coverage radius $\rho_{\mathrm{cov}} \in (0,\rho_{\mathrm{bind}}]$, where $\rho_{\mathrm{bind}}$ is
the binding radius associated with the affinity threshold
$\Delta G_{\mathrm{thresh}}$.
A footprint $V \subseteq \mathcal{I}^R$ \emph{strongly covers} an antigen path
$\Gamma_A \in \mathcal{P}_A(L_A^\star)$ at radius $\rho_{\mathrm{cov}}$ if
\begin{equation}
  d_{\max}(\Gamma_A,V) \;\le\; \rho_{\mathrm{cov}}.
\end{equation}
Equivalently, every antigenic state along $\Gamma_A$ lies within mirror-space distance
$\rho_{\mathrm{cov}}$ of at least one anchor $v \in V$.
\end{definition}

This is a purely geometric notion: it depends only on $(\mathcal{M}_A,\mathcal{M}_R)$, the
co-embedding, and the path family. In the probabilistic setting, with a path distribution
$\Pi_A(L_A^\star)$, we can define a strong-coverage functional
\begin{equation}
  C_{\mathrm{cov}}^{\mathrm{strong}}(V;\rho_{\mathrm{cov}})
  := \Pr_{\Gamma_A \sim \Pi_A(L_A^\star)}
     \bigl( d_{\max}(\Gamma_A,V) \le \rho_{\mathrm{cov}} \bigr).
\end{equation}

\paragraph{Worst-case uncovered geodesic distance.}
In many settings it is also natural to measure the worst-case geodesic distance along a path to
the nearest vaccine anchor. For $z \in \mathcal{M}_A$ and $V \subseteq \mathcal{M}_R$ define
\[
  d_{g_A}(z,V) := \inf_{v \in V} d_{g_A}\bigl(z,v\bigr),
\]
where we regard anchors $v$ as living in a common mirror space via the complementary
co-embedding. A worst-case uncovered distance functional is
\begin{equation}
  U_{\mathrm{geo}}(V)
  := \sup_{\Gamma_A \in \mathcal{P}_A(L_A^\star)}
     \;\sup_{t \in [0,1]} d_{g_A}\bigl(\Gamma_A(t),V\bigr).
\end{equation}
Minimizing $U_{\mathrm{geo}}(V)$ seeks vaccine footprints that keep all plausible drift trajectories
geodesically close to at least one induced antibody cluster.

\paragraph{Cantelli-based success lower bounds.}
The geometric notions above connect to neutralization via the binding calibration and inverted
error budget. Given a vaccine footprint $V \subseteq \mathcal{I}^R$, the MIRADOR pipeline
associates a \emph{neutralization coverage certificate}
\begin{equation}
  \underline{C}_{\mathrm{neut}}(V;L_A^\star)
  := \bigl[\widehat{C}(V;L_A^\star) - \varepsilon_{\mathrm{samp}}\bigr]_+
     - \bigl( \overline{E}_{\mathrm{geom}} + \overline{E}_{\mathrm{bind}}
              + \overline{E}_{\mathrm{link}} \bigr),
\end{equation}
where $\widehat{C}$ is the empirical binding-coverage estimate over Monte Carlo antigen paths at
horizon $L_A^\star$, $\varepsilon_{\mathrm{samp}}$ is a sampling tolerance (e.g., from a Hoeffding
bound), and $(\overline{E}_{\mathrm{geom}},\overline{E}_{\mathrm{bind}},\overline{E}_{\mathrm{link}})$ are
conservative upper bounds on the geometric, binding, and linkage error-budget terms. The
certificate $\underline{C}_{\mathrm{neut}}(V;L_A^\star)$ is a Cantelli-based lower bound on the
true neutralization coverage $C_{\mathrm{neut}}(V;L_A^\star)$.

\paragraph{Vaccine design as a constrained coverage problem.}
Let $\mathcal{V}$ denote a class of feasible vaccine footprints (e.g., all $K$-anchor footprints
induced by realizable immunogens under manufacturing and regulatory constraints), and let
$C_\star \in (0,1)$ be a target neutralization coverage. MIRADOR expresses vaccine design as
\begin{equation}
\label{eq:mirador-vaccine-design}
  \text{find } V \in \mathcal{V}
  \quad\text{such that}\quad
  \underline{C}_{\mathrm{neut}}(V;L_A^\star) \;\ge\; C_\star,
\end{equation}
optionally with an additional geometric constraint $U_{\mathrm{geo}}(V) \le R_{\max}$ for a
pre-registered geodesic radius $R_{\max}$. In practice, one can either:
\begin{itemize}
  \item \emph{minimize} $U_{\mathrm{geo}}(V)$ over $\mathcal{V}$ subject to
        $\underline{C}_{\mathrm{neut}}(V;L_A^\star) \ge C_\star$, or
  \item \emph{maximize} $\underline{C}_{\mathrm{neut}}(V;L_A^\star)$ subject to
        $U_{\mathrm{geo}}(V) \le R_{\max}$ and panel-size or manufacturability constraints.
\end{itemize}
Equation~\eqref{eq:mirador-vaccine-design} turns the abstract coverage certificate into a concrete
design requirement.

\subsection{Greedy and relaxed optimization in mirror space}
\label{subsec:greedy-relaxed}

Solving~\eqref{eq:mirador-vaccine-design} exactly is combinatorial: even with a finite library of
candidate immunogens, the number of possible footprints $V$ grows rapidly with $K$. MIRADOR
therefore relies on approximate algorithms that operate directly in mirror space.

\paragraph{Greedy mirror-space covering.}
Suppose we have a discrete candidate set
$\mathcal{U} \subseteq \mathcal{I}^R$ of receptor-like anchors induced by a library of immunogens
(e.g., historical strains, prototype constructs, or outputs of a generative model). Define a set
function
\[
  F(V) := \underline{C}_{\mathrm{neut}}(V;L_A^\star),
  \qquad V \subseteq \mathcal{U},
\]
using a Monte Carlo estimate of $\widehat{C}(V;L_A^\star)$ and a fixed error budget. Because adding
anchors cannot reduce coverage, $F$ is monotone, and in many regimes it behaves approximately
submodular.

A simple greedy algorithm initializes $V_0 = \varnothing$ and iteratively chooses
\[
  v_{m} \in \arg\max_{v \in \mathcal{U} \setminus V_{m-1}}
               \bigl( F(V_{m-1} \cup \{v\}) - F(V_{m-1}) \bigr),
\]
until a panel-size budget is reached or the marginal gain falls below a tolerance. The output
$V_M$ is then plugged into the coverage certificate to obtain a MIRADOR-ready vaccine design.

\paragraph{Continuous relaxations and projection back to sequence space.}
In settings where the design space is not easily enumerated, we can relax the problem by
optimizing over anchors $z_1,\dots,z_K \in \mathcal{M}_R$ directly:
\begin{equation}
\label{eq:relaxed-objective}
  J(z_1,\dots,z_K)
  := \mathbb{E}_{\Gamma_A \sim \Pi_A(L_A^\star)}
      \bigl[ \sigma\bigl(-d_{\max}(\Gamma_A,\{z_1,\dots,z_K\})\bigr) \bigr],
\end{equation}
where $\sigma$ is a smooth, monotonically increasing surrogate (e.g., logistic). Riemannian
gradient ascent on the product manifold $\mathcal{M}_R^K$ yields a set of \emph{ideal} anchors
$\{z_k^\star\}$ that maximize the surrogate coverage functional.

To obtain realizable immunogens, we then apply a \emph{projection} operator
$\Pi_{\mathrm{seq}} : \mathcal{M}_R \to \text{Seq}$ that maps each ideal anchor $z_k^\star$ to a
sequence or construct whose receptor footprint $u_R(\mathcal{R}_{\nu_k})$ lies near $z_k^\star$
while respecting biophysical and manufacturing constraints. This projection may be implemented
via nearest-neighbor search in a large design database, via conditional generative models, or via
local sequence optimization constrained to stay within a small neighborhood of $z_k^\star$.

The resulting footprint
\[
  V_{\mathrm{relaxed}} := \Pi_{\mathrm{seq}}(z_1^\star) \cup \dots \cup \Pi_{\mathrm{seq}}(z_K^\star)
\]
is then evaluated with the full MIRADOR certificate
$\underline{C}_{\mathrm{neut}}(V_{\mathrm{relaxed}};L_A^\star)$ and either accepted (if it meets
$C_\star$) or refined via additional rounds of optimization.

\subsection{Antibody discovery and optimization in mirror space}
\label{subsec:antibody-design}

MIRADOR also supports design on the receptor side: given a family of antigens (current and
plausible future states), we can search for therapeutic antibodies or receptor panels that maximize
certified coverage.

\paragraph{Single-antibody optimization.}
Let $\mathcal{A}_{\mathrm{focus}} \subseteq \mathcal{A}$ be a set of antigens of interest and
$\Pi_A(L_A^\star)$ a path distribution supported on their future drift. For a single candidate
receptor $r$ we can define a success functional
\begin{equation}
  J_{\mathrm{single}}(r)
  := \underline{C}_{\mathrm{neut}}(\{r\};L_A^\star),
\end{equation}
where the certificate is computed with $V = \{r\}$ and includes both geometric and binding
error-budget terms. MIRADOR seeks
\[
  r^\star \in \arg\max_{r \in \mathcal{R}_{\mathrm{design}}} J_{\mathrm{single}}(r),
\]
where $\mathcal{R}_{\mathrm{design}}$ encodes developability and liability constraints via the
receptor configuration map and margins (e.g., productive, properly folded, non-autoreactive).

In mirror space this becomes a continuous optimization problem over $z_R = u_R(r) \in \mathcal{M}_R$:
one can approximate the gradient $\nabla_{z_R} J_{\mathrm{single}}$ by Monte Carlo over antigen
paths and perform Riemannian gradient ascent, followed by projection back to sequence space via
an antibody generative model or structure-aware sequence search.

\paragraph{Panel design and benign receptor paths.}
For therapeutic cocktails we consider panels $W = \{r_1,\dots,r_M\} \subseteq \mathcal{I}^R$ and
optimize $J_{\mathrm{panel}}(W) := \underline{C}_{\mathrm{neut}}(W;L_A^\star)$ using the same
greedy or relaxed strategies as for vaccine footprints. An additional degree of freedom is the
affinity-maturation trajectory of each antibody: benign receptor paths
$\mathcal{P}_R(L_R^\star)$ in $\mathcal{M}_R$ describe how a germline precursor can evolve under
somatic hypermutation and selection.

\paragraph{Mirror-aware affinity maturation.}
Because MIRADOR provides a Riemannian geometry on the receptor side, it suggests
\emph{mirror-aware} affinity-maturation protocols that follow high-affinity geodesics. Starting from
a germline seed $r_{\mathrm{germ}}$ with embedding $z_{\mathrm{germ}} = u_R(r_{\mathrm{germ}})$,
one can:
\begin{enumerate}
  \item define a surrogate objective $J_{\mathrm{mat}}(z_R)$ that rewards strong binding along
        antigen paths (e.g., expected minimum distance to $\Gamma_A(t)$ over $\Pi_A$);
  \item perform geodesic or gradient-based updates within $\mathcal{P}_R(L_R^\star)$,
        \[
          z_R^{(m+1)} = \exp_{z_R^{(m)}}\bigl( \eta \,\mathrm{grad}_{z_R} J_{\mathrm{mat}}(z_R^{(m)}) \bigr),
        \]
        with step size $\eta > 0$ and retraction via the Riemannian exponential; and
  \item map intermediate points $z_R^{(m)}$ back to antibody sequences via structure- or
        sequence-based generative models.
\end{enumerate}
This yields \emph{in vitro} or \emph{in silico} maturation trajectories that explicitly track
high-affinity regions of mirror space rather than operating blindly in sequence space.

\subsection{Joint design of vaccines and therapeutic cocktails}
\label{subsec:joint-design}

In many realistic programs vaccines and therapeutic antibodies are deployed together. MIRADOR’s
mirror geometry naturally supports \emph{joint} design of vaccine footprints and therapeutic panels.

Let $V_{\mathrm{vac}} \subseteq \mathcal{I}^R$ denote the footprint induced by a vaccine strategy
(e.g., a multivalent mRNA payload plus schedule), and let
$V_{\mathrm{tx}} \subseteq \mathcal{I}^R$ denote a therapeutic cocktail. For an antigen path
$\Gamma_A$ define the neutralization hit events
\[
  \Hit_{\mathrm{vac}}(\Gamma_A) := \Hit_{\mathrm{neut}}(V_{\mathrm{vac}},\Gamma_A),
  \qquad
  \Hit_{\mathrm{tx}}(\Gamma_A)  := \Hit_{\mathrm{neut}}(V_{\mathrm{tx}},\Gamma_A).
\]
The joint success event is
\[
  \Hit_{\mathrm{joint}}(\Gamma_A)
  := \Hit_{\mathrm{vac}}(\Gamma_A) \,\vee\, \Hit_{\mathrm{tx}}(\Gamma_A),
\]
so that at least one of the two defenses neutralizes the path. The corresponding joint coverage is
\begin{equation}
  C_{\mathrm{joint}}(V_{\mathrm{vac}},V_{\mathrm{tx}};L_A^\star)
  := \Pr_{\Gamma_A \sim \Pi_A(L_A^\star)}
       \bigl( \Hit_{\mathrm{joint}}(\Gamma_A) \bigr).
\end{equation}
Using inclusion–exclusion,
\[
  C_{\mathrm{joint}}
  = C_{\mathrm{neut}}(V_{\mathrm{vac}};L_A^\star)
    + C_{\mathrm{neut}}(V_{\mathrm{tx}};L_A^\star)
    - \Pr\bigl( \Hit_{\mathrm{vac}} \wedge \Hit_{\mathrm{tx}} \bigr),
\]
and analogous expressions can be written for joint coverage certificates
$\underline{C}_{\mathrm{joint}}(V_{\mathrm{vac}},V_{\mathrm{tx}};L_A^\star)$ with an additional
independence-slack term capturing correlations between the two defenses.

A joint design problem then seeks
\begin{equation}
\label{eq:joint-design}
  (V_{\mathrm{vac}}^\star,V_{\mathrm{tx}}^\star)
  \in \arg\max_{V_{\mathrm{vac}} \in \mathcal{V}_{\mathrm{vac}},
               \, V_{\mathrm{tx}} \in \mathcal{V}_{\mathrm{tx}}}
      \underline{C}_{\mathrm{joint}}(V_{\mathrm{vac}},V_{\mathrm{tx}};L_A^\star),
\end{equation}
subject to separate constraints on vaccine manufacturability, dosing, and therapeutic
developability. Geometrically, this encourages vaccine footprints and therapeutic panels to cover
complementary regions of mirror space, minimizing the probability that an antigenic path lies
outside the coverage of both.

\subsection{Toy scenarios (conceptual)}
\label{subsec:toy-scenarios}

The following thought experiments sketch how MIRADOR could be applied in concrete pathogen
settings; they are conceptual only and do not report empirical results.

\paragraph{SARS-CoV-2: hedging against HERALD-predicted drift.}
Suppose HERALD has been instantiated on SARS-CoV-2 and provides a path distribution
$\Pi_A^{\mathrm{S2}}(L_A^\star)$ that captures likely antigenic drift over the next $1$–$2$ years,
with paths concentrated along RBD-escape and NTD-escape directions. MIRADOR receives this
distribution and a library $\mathcal{U}_{\mathrm{S2}}$ of candidate spike immunogens (including
existing and prototype mRNA constructs).

Two design options might be:
\begin{itemize}
  \item \emph{Monovalent booster.} Choose a single immunogen $\nu_{\mathrm{mono}}$ whose footprint
        $V_{\mathrm{mono}}$ is closest (in geodesic distance) to the current dominant variant.
        This yields very high coverage for short horizons but may leave some HERALD-predicted drift
        trajectories poorly covered, resulting in a modest
        $\underline{C}_{\mathrm{neut}}(V_{\mathrm{mono}};L_A^\star)$.
  \item \emph{Hedged multivalent payload.} Use greedy or relaxed mirror-space optimization to
        select two or three immunogens $\nu_1,\nu_2,\nu_3$ with footprints
        $V_{\mathrm{hedge}} = V_{\nu_1} \cup V_{\nu_2} \cup V_{\nu_3}$ that flank the main drift
        directions in mirror space. This may slightly reduce near-term coverage but increase the
        certified coverage $\underline{C}_{\mathrm{neut}}(V_{\mathrm{hedge}};L_A^\star)$ across the
        entire HERALD ensemble.
\end{itemize}
MIRADOR makes this trade-off explicit: the design choice is not between “good” and “bad”
vaccines, but between different shapes of coverage certificates over the path ensemble.

\paragraph{Influenza: universal constructs versus strain-matched vaccines.}
Consider an H3N2 manifold where historical antigenic drift traces long, roughly linear paths with
occasional branch jumps. A “universal” construct (e.g., stem-focused HA) induces a footprint
$V_{\mathrm{univ}}$ near the center of $\mathcal{M}_A$, while strain-matched vaccines induce
footprints $V_{\mathrm{SM}}(t)$ that chase the currently dominant variant each season.

In MIRADOR, one could:
\begin{itemize}
  \item estimate $\Pi_A^{\mathrm{flu}}(L_A^\star)$ from historical and model-based drift;
  \item compute coverage certificates
        $\underline{C}_{\mathrm{neut}}(V_{\mathrm{univ}};L_A^\star)$ and
        $\underline{C}_{\mathrm{neut}}(V_{\mathrm{SM}};L_A^\star)$ under different update cadences; and
  \item evaluate mixed strategies (e.g., combining a universal primer with periodic strain-matched
        boosts) as joint designs of the form~\eqref{eq:joint-design}.
\end{itemize}
Even without precise numerical values, MIRADOR frames the question geometrically: universal
constructs may achieve moderate but stable coverage certificates over a broad set of paths, whereas
purely strain-matched strategies may achieve high short-term coverage but rapidly degrade as
$\Pi_A^{\mathrm{flu}}$ explores new directions. The mirror geometry and coverage certificates make
these trade-offs analyzable rather than purely qualitative.


\section{Inverted error budgets and coverage certificates}
\label{sec:inverted-error-budgets}

MIRADOR's construction gives us a mirror geometry and a complementary co-embedding in which
binding events are modeled by Euclidean distances between antigen and receptor embeddings
(Definition~\ref{def:complementary-coembedding}). In this section we \emph{invert} the Davis
error-budget viewpoint: instead of asking how often a detector is correct when it fires, we ask
how often a receptor panel $V$ neutralizes forecasted antigenic drift paths, and we derive
\emph{coverage certificates}---lower bounds on neutralization coverage that explicitly account for
geometry, binding, linkage, and sampling error.

Conceptually, we mirror the Davis error-budget decomposition for detection systems, but with
escape/detection events replaced by path-wise interception of antigenic drift trajectories by
designed receptor panels.\footnote{For the original detection-centric error-budget formulation,
see the Davis manifold framework and HERALD's dominance bounds.}

Throughout this section we fix a path horizon $L \le L_A^\star$ and a forecast distribution
$\Pi_A(L)$ over benign antigen paths $\Gamma_A : [0,1] \to \mathcal{M}_A$ (Definition~\ref{def:mirror-paths}).
As in Section~\ref{sec:mirror-geometry}, we write
\[
\delta\bigl(\Gamma_A(t),r\bigr)
\quad\text{as shorthand for}\quad
\delta\bigl(u_A(\Gamma_A(t)),u_R(r)\bigr),
\]
and use the positive part notation $[u]_+ := \max\{u,0\}$.

\subsection{Three levels of coverage and hit events}

We begin by distinguishing three levels of coverage: true neutralization coverage, model-level
binding coverage, and empirical Monte Carlo coverage.

\begin{definition}[Hit events and coverage levels]
\label{def:coverage-levels}
Fix a horizon $L \le L_A^\star$, a receptor panel $V \subseteq \mathcal{R}$, and an antigenic
drift path $\Gamma_A : [0,1] \to \mathcal{M}_A$ drawn from $\Pi_A(L)$.

\begin{enumerate}
\item \emph{Binding hit event.} The binding hit event is
\[
\Hit_{\mathrm{bind}}(V,\Gamma_A)
\;:=\;
\Bigl\{
\exists\,t \in [0,1],\ \exists\,r \in V :
\delta\bigl(\Gamma_A(t),r\bigr) \le \rho_{\mathrm{bind}}
\Bigr\},
\]
where $\rho_{\mathrm{bind}} = \Phi_{\mathrm{bind}}^{-1}(\Delta G_{\mathrm{thresh}})$ is the binding-distance
threshold from Definition~\ref{def:complementary-coembedding}.

\item \emph{Neutralization hit event.} The neutralization hit event
$\Hit_{\mathrm{neut}}(V,\Gamma_A)$ is the ground-truth event that along the path $\Gamma_A$ there
exists a time $t$ at which at least one $r \in V$ neutralizes the corresponding antigenic state
at or above a pre-specified potency threshold (e.g., IC$_{50}$ or fold-drop criterion). This event
is defined in the underlying biological process or simulator, not by the embedding.

\item \emph{True neutralization coverage.} The \emph{true} coverage of $V$ at horizon $L$ is
\[
C_{\mathrm{neut}}(V;L)
\;:=\;
\Pr_{\Gamma_A \sim \Pi_A(L)}
\Bigl(\Hit_{\mathrm{neut}}(V,\Gamma_A)\Bigr).
\]

\item \emph{Model binding coverage.} The \emph{model-level} binding coverage of $V$ at horizon
$L$ is
\[
C_{\mathrm{bind}}(V;L)
\;:=\;
\Pr_{\Gamma_A \sim \Pi_A(L)}
\Bigl(\Hit_{\mathrm{bind}}(V,\Gamma_A)\Bigr).
\]
When the horizon is clear from context we write $C(V;L)$ or simply $C(V)$ for
$C_{\mathrm{bind}}(V;L)$.
\end{enumerate}
\end{definition}

The rest of the section studies how to lower-bound $C_{\mathrm{neut}}(V;L)$ in terms of
$C_{\mathrm{bind}}(V;L)$ and an empirical estimator of the latter, with an explicit coverage
error budget.

\subsection{Empirical coverage and sampling tolerance}

In practice, $C_{\mathrm{bind}}(V;L)$ is estimated by Monte Carlo over a finite sample of
forecasted drift paths.

\begin{definition}[Empirical binding coverage and sampling tolerance]
\label{def:empirical-coverage}
Let $\Gamma_A^{(1)},\dots,\Gamma_A^{(B)}$ be i.i.d.\ draws from $\Pi_A(L)$, and define the
empirical binding coverage of $V$ at horizon $L$ as
\[
\widehat{C}(V;L)
\;:=\;
\frac{1}{B} \sum_{b=1}^B
\mathbf{1}\!\Bigl\{
\Hit_{\mathrm{bind}}\bigl(V,\Gamma_A^{(b)}\bigr)
\Bigr\}.
\]

For a target sampling confidence level $\delta_{\mathrm{samp}} \in (0,1)$, we define the
\emph{sampling tolerance}
\begin{equation}
\label{eq:sampling-tolerance}
\varepsilon_{\mathrm{samp}}
\;:=\;
\sqrt{\frac{1}{2B}\,
\log\!\Bigl(\frac{2}{\delta_{\mathrm{samp}}}\Bigr)}.
\end{equation}
By Hoeffding's inequality, with probability at least $1-\delta_{\mathrm{samp}}$ we have
\[
\bigl|\widehat{C}(V;L) - C_{\mathrm{bind}}(V;L)\bigr|
\;\le\;
\varepsilon_{\mathrm{samp}}.
\]
\end{definition}

Thus $\widehat{C}(V;L)$ is an unbiased estimator of $C_{\mathrm{bind}}(V;L)$ with a finite-sample
confidence band controlled by $\varepsilon_{\mathrm{samp}}$ and $B$.

\subsection{Coverage error budget and independence slack}

We now decompose the discrepancy between $C_{\mathrm{neut}}(V;L)$ and $C_{\mathrm{bind}}(V;L)$ into
three interpretable model components, and introduce an independence slack that captures
interactions between them.

For a random path $\Gamma_A \sim \Pi_A(L)$, consider the following ``good'' events:

\begin{itemize}
\item $G_{\mathrm{geom}}$: the antigen path remains in the validated geometric regime on
$\mathcal{M}_A$ and within the binding in-distribution region, so that bounded-distortion and
margin assumptions hold along $\Gamma_A$;

\item $G_{\mathrm{bind}}$: conditional on $G_{\mathrm{geom}}$, the binding surrogate
$d(\Gamma_A(t),r) \mapsto \Delta G(a(t),r)$ correctly classifies binding vs.~non-binding
relative to $\Delta G_{\mathrm{thresh}}$ for all $t$ and $r \in V$ of interest;

\item $G_{\mathrm{link}}$: conditional on $G_{\mathrm{geom}} \cap G_{\mathrm{bind}}$, the linkage between
binding and neutralization holds: if the binding event fires for some $(t,r)$, then so does the
neutralization event at that point, and vice versa, up to rare linkage failures.
\end{itemize}

Write $G_{\mathrm{cov}} := G_{\mathrm{geom}} \cap G_{\mathrm{bind}} \cap G_{\mathrm{link}}$ and define the
marginal coverage error probabilities
\[
E_{\mathrm{geom}} := \Pr\bigl(G_{\mathrm{geom}}^{\mathrm{c}}\bigr),\qquad
E_{\mathrm{bind}} := \Pr\bigl(G_{\mathrm{bind}}^{\mathrm{c}}\bigr),\qquad
E_{\mathrm{link}} := \Pr\bigl(G_{\mathrm{link}}^{\mathrm{c}}\bigr).
\]

\begin{definition}[Coverage independence slack]
\label{def:coverage-independence}
The \emph{coverage independence slack} $\delta_{\mathrm{indep}}^{\mathrm{cov}} \in [0,1]$ is
\[
\delta_{\mathrm{indep}}^{\mathrm{cov}}
\;:=\;
\Bigl[
(1 - E_{\mathrm{geom}})(1 - E_{\mathrm{bind}})(1 - E_{\mathrm{link}})
- \Pr\bigl(G_{\mathrm{cov}}\bigr)
\Bigr]_+,
\]
where $[u]_+ := \max\{u,0\}$ denotes the positive part.
\end{definition}

If $G_{\mathrm{geom}}, G_{\mathrm{bind}}, G_{\mathrm{link}}$ are independent, then
$\Pr(G_{\mathrm{cov}}) = \prod (1-E_{\bullet})$ and
$\delta_{\mathrm{indep}}^{\mathrm{cov}} = 0$. When these events tend to fail together, the actual
good-event probability $\Pr(G_{\mathrm{cov}})$ may be smaller than the product, and
$\delta_{\mathrm{indep}}^{\mathrm{cov}}$ measures this shortfall.

\subsection{Compositional coverage bounds}

We can now state a compositional bound that lower-bounds true neutralization coverage in terms
of empirical binding coverage and the coverage error budget.

\begin{theorem}[Compositional coverage bound]
\label{thm:coverage-bound}
Fix a horizon $L \le L_A^\star$ and a panel $V \subseteq \mathcal{R}$. Let
$C_{\mathrm{neut}}(V;L)$ and $C_{\mathrm{bind}}(V;L)$ be as in
Definition~\ref{def:coverage-levels}, and let $\widehat{C}(V;L)$ and
$\varepsilon_{\mathrm{samp}}$ be as in Definition~\ref{def:empirical-coverage}.

Assume:
\begin{enumerate}
\item (\emph{Geometric validity}) The Davis mirror-geometry assumptions hold on the
forecast distribution $\Pi_A(L)$ with failure probability at most $E_{\mathrm{geom}}$; that is,
$\Pr(G_{\mathrm{geom}}^{\mathrm{c}}) \le E_{\mathrm{geom}}$.

\item (\emph{Binding-model validity}) Conditional on $G_{\mathrm{geom}}$, the binding surrogate
correctly classifies binding above/below $\Delta G_{\mathrm{thresh}}$ with failure probability at most
$E_{\mathrm{bind}}$.

\item (\emph{Binding--neutralization linkage}) Conditional on
$G_{\mathrm{geom}} \cap G_{\mathrm{bind}}$, the neutralization hit event agrees with the binding hit
event with failure probability at most $E_{\mathrm{link}}$.

\item (\emph{Sampling}) The empirical coverage is computed from $B$ i.i.d.\ draws from
$\Pi_A(L)$, and $\varepsilon_{\mathrm{samp}}$ is chosen according to
\eqref{eq:sampling-tolerance}, so that with probability at least $1-\delta_{\mathrm{samp}}$,
\[
\bigl|\widehat{C}(V;L) - C_{\mathrm{bind}}(V;L)\bigr| \le \varepsilon_{\mathrm{samp}}.
\]
\end{enumerate}

Then, with probability at least $1-\delta_{\mathrm{samp}}$ over the Monte Carlo sampling,
\begin{equation}
\label{eq:coverage-additive}
C_{\mathrm{neut}}(V;L)
\;\ge\;
\bigl[\widehat{C}(V;L) - \varepsilon_{\mathrm{samp}}\bigr]_+
\;-\;
\bigl(E_{\mathrm{geom}} + E_{\mathrm{bind}} + E_{\mathrm{link}}\bigr),
\end{equation}
and, moreover,
\begin{equation}
\label{eq:coverage-multiplicative}
C_{\mathrm{neut}}(V;L)
\;\ge\;
\bigl[\widehat{C}(V;L) - \varepsilon_{\mathrm{samp}}\bigr]_+
\Bigl(
(1 - E_{\mathrm{geom}})(1 - E_{\mathrm{bind}})(1 - E_{\mathrm{link}})
- \delta_{\mathrm{indep}}^{\mathrm{cov}}
\Bigr),
\end{equation}
where $\delta_{\mathrm{indep}}^{\mathrm{cov}}$ is the coverage independence slack from
Definition~\ref{def:coverage-independence}. In regimes where the right-hand side of
\eqref{eq:coverage-multiplicative} is negative, it is clipped at zero.
\end{theorem}

\begin{proof}[Proof sketch]
On the coverage-good event $G_{\mathrm{cov}}$, the binding and neutralization hit events coincide
along each path, so $C_{\mathrm{neut}}(V;L)$ and $C_{\mathrm{bind}}(V;L)$ differ only on
$G_{\mathrm{cov}}^{\mathrm{c}}$. A union bound over the three failure modes yields
\[
C_{\mathrm{neut}}(V;L)
\;\ge\;
C_{\mathrm{bind}}(V;L) - \bigl(E_{\mathrm{geom}} + E_{\mathrm{bind}} + E_{\mathrm{link}}\bigr).
\]
Substituting the sampling inequality
$C_{\mathrm{bind}}(V;L) \ge \widehat{C}(V;L) - \varepsilon_{\mathrm{samp}}$ and truncating at zero
gives~\eqref{eq:coverage-additive}.

For the multiplicative form, note that
$\Pr(G_{\mathrm{cov}}) \ge (1 - E_{\mathrm{geom}})(1 - E_{\mathrm{bind}})(1 - E_{\mathrm{link}})
- \delta_{\mathrm{indep}}^{\mathrm{cov}}$ by
Definition~\ref{def:coverage-independence}. Conditioning on $G_{\mathrm{cov}}$ and using that
$\Hit_{\mathrm{neut}} = \Hit_{\mathrm{bind}}$ on this event, we obtain
\[
C_{\mathrm{neut}}(V;L)
\;\ge\;
\Pr(G_{\mathrm{cov}})
\,
C_{\mathrm{bind}}(V;L),
\]
which combined with the sampling bound yields~\eqref{eq:coverage-multiplicative} after clipping
negative values. A full proof simply spells out these steps and tracks constants; the arguments
are elementary.
\end{proof}

Conceptually, \eqref{eq:coverage-additive} is a conservative, assumptions-light lower bound, while
\eqref{eq:coverage-multiplicative} exploits approximate independence between coverage error
sources when justified by validation.

\subsection{Coverage certificates and inverted budgets}

Theorems of the form~\eqref{eq:coverage-additive}--\eqref{eq:coverage-multiplicative} naturally
lead to \emph{coverage certificates}: objects that certify that a given panel achieves at least a
target neutralization coverage under the stated assumptions.

\begin{definition}[Coverage certificate]
\label{def:coverage-certificate}
Fix a target coverage level $C_\star \in (0,1)$, a horizon $L \le L_A^\star$, and confidence
parameters $(\delta_{\mathrm{samp}},\delta_{\mathrm{model}})$. A
$(1-\delta_{\mathrm{samp}}-\delta_{\mathrm{model}})$-\emph{coverage certificate} for a panel
$V \subseteq \mathcal{R}$ at horizon $L$ consists of:
\begin{itemize}
\item an empirical coverage estimate $\widehat{C}(V;L)$ computed from $B$ Monte Carlo paths;
\item a sampling tolerance $\varepsilon_{\mathrm{samp}}$ as in \eqref{eq:sampling-tolerance};
\item upper bounds $\overline{E}_{\mathrm{geom}}, \overline{E}_{\mathrm{bind}},
\overline{E}_{\mathrm{link}}$ on the coverage error components and
$\overline{\delta}_{\mathrm{indep}}^{\mathrm{cov}}$ on the independence slack, each validated on
held-out data with probability at least $1-\delta_{\mathrm{model}}$; and
\item a verification that, after plugging these quantities into either
\eqref{eq:coverage-additive} or \eqref{eq:coverage-multiplicative}, the resulting lower bound on
$C_{\mathrm{neut}}(V;L)$ is at least $C_\star$.
\end{itemize}
We say that $V$ is \emph{$C_\star$-certified at horizon $L$} when such a certificate exists.
\end{definition}

From the additive bound~\eqref{eq:coverage-additive} we can read off a simple inverted form: any
panel $V$ whose empirical coverage exceeds a threshold
\begin{equation}
\label{eq:required-empirical-coverage}
\widehat{C}_{\min}(C_\star)
\;:=\;
C_\star
+ \varepsilon_{\mathrm{samp}}
+ \overline{E}_{\mathrm{geom}}
+ \overline{E}_{\mathrm{bind}}
+ \overline{E}_{\mathrm{link}}
\end{equation}
necessarily satisfies $C_{\mathrm{neut}}(V;L) \ge C_\star$ under the corresponding error bounds.
In this sense, the coverage error budget is \emph{inverted}: instead of propagating modeling
errors forward to obtain a lower bound, we use the budget to specify how large
$\widehat{C}(V;L)$ must be for a design to be certifiable.

\begin{proposition}[Inverted coverage budget]
\label{prop:inverted-budget}
Let $C_\star \in (0,1)$ and suppose that for a panel $V$ and horizon $L$ we have empirical
coverage $\widehat{C}(V;L)$, sampling tolerance $\varepsilon_{\mathrm{samp}}$, and upper bounds
$\overline{E}_{\mathrm{geom}}, \overline{E}_{\mathrm{bind}}, \overline{E}_{\mathrm{link}}$ satisfying
\[
\widehat{C}(V;L)
\;\ge\;
\widehat{C}_{\min}(C_\star)
\]
with $\widehat{C}_{\min}(C_\star)$ as in \eqref{eq:required-empirical-coverage}. Then, under the
assumptions of Theorem~\ref{thm:coverage-bound}, $C_{\mathrm{neut}}(V;L) \ge C_\star$ with
probability at least $1-\delta_{\mathrm{samp}}-\delta_{\mathrm{model}}$. Conversely, if
$\widehat{C}(V;L) < \widehat{C}_{\min}(C_\star)$, the additive bound~\eqref{eq:coverage-additive}
cannot certify $C_{\mathrm{neut}}(V;L) \ge C_\star$.
\end{proposition}

In the multiplicative setting~\eqref{eq:coverage-multiplicative}, an analogous inverted expression
involves dividing by
$(1 - \overline{E}_{\mathrm{geom}})(1 - \overline{E}_{\mathrm{bind}})(1 - \overline{E}_{\mathrm{link}})
- \overline{\delta}_{\mathrm{indep}}^{\mathrm{cov}}$, provided this factor is positive and not too
small. In regimes where this factor is close to zero, the multiplicative bound becomes numerically
fragile and the additive form is preferable.

\subsection{Practical use, vacuity, and slice-aware certificates}

In practice, coverage certificates are most informative when they are:
\begin{enumerate}
\item \emph{slice-aware}: defined not only for a global path distribution $\Pi_A(L)$, but also
for epidemiologically meaningful slices (e.g., drift paths starting from a particular vaccine
anchor, restricted to a lineage, or constrained to an epitope subset);

\item \emph{explicitly non-vacuous}: accompanied by a check that the right-hand side of
\eqref{eq:coverage-additive} or \eqref{eq:coverage-multiplicative} is appreciably above zero
for the slice in question; and

\item \emph{paired with abstention}: when the bound is near zero or negative, MIRADOR should
treat it as vacuous and either abstain from issuing a certificate for that slice or tighten the
scope (e.g., shorter horizons, smaller path families).
\end{enumerate}

Operationally, MIRADOR's design goals are to construct panels $V$ that:
\begin{enumerate}
\item \emph{achieve high coverage} $C(V;L)$ against drift families of interest; and
\item \emph{provide a coverage certificate}: a lower bound on $C_{\mathrm{neut}}(V;L)$ that folds
in geometric distortion, binding-model uncertainty, linkage error between binding and
neutralization, and sampling error.
\end{enumerate}
The inverted error-budget view makes these goals precise: the budget specifies ``how much slack''
must be added to a desired coverage level $C_\star$ to obtain the required empirical coverage
$\widehat{C}_{\min}(C_\star)$, and MIRADOR's optimization problem is to design panels that clear
this bar for as many relevant slices and horizons as possible.

\section{Ethics, Dual-Use, and Generative Firewalls}
\label{sec:ethics-dual-use}

MIRADOR differs from HERALD in a critical respect: whereas HERALD performs retrospective
surveillance and detection, MIRADOR prospectively \emph{designs} immunological objects (vaccine
footprints, antibody panels, receptor sets). This shift from passive observation to active design
introduces dual-use risks that require explicit architectural and governance safeguards. This
section articulates MIRADOR's threat model, specifies banned optimization modes, and documents
the information-hazard mitigations that must be enforced in any deployment.

\subsection{Dual-use threat model for MIRADOR}
\label{subsec:threat-model}

The central dual-use risk is \emph{inverse inference}: if MIRADOR publicly releases an
``optimal'' receptor or panel $V^\star$ that maximally covers a forecasted drift space, an
adversary may invert the complementary co-embedding to infer which future antigenic states
$a^\star \in \mathcal{M}_A$ are \emph{least} covered by $V^\star$, effectively obtaining a
blueprint for immune-evading or synthetically optimized pathogens.

\paragraph{Threat scenarios.}
\begin{enumerate}
  \item \textbf{Receptor-to-antigen inversion.}
  Given a highly optimized receptor $r^\star$ or panel $V^\star$, an adversary with access to
  the mirror geometry $(\mathcal{M}_A, \mathcal{M}_R, u_A, u_R)$ solves
  \[
    a_{\mathrm{evade}} \in \arg\min_{a \in \mathcal{M}_A} \;\min_{v \in V^\star}
                       \delta\bigl(u_A(a), u_R(v)\bigr),
  \]
  yielding an antigen that maximally evades $V^\star$. If $V^\star$ is globally optimal against
  natural drift, $a_{\mathrm{evade}}$ may represent a worst-case synthetic threat.

  \item \textbf{Path-family amplification.}
  MIRADOR's path distribution $\Pi_A(L_A^\star)$ is derived from surveillance and mechanistic
  models. If this distribution or the underlying benign path family $\mathcal{P}_A(L_A^\star)$
  is published in full detail, adversaries may use it to identify ``natural-looking'' drift
  trajectories with maximal immune-escape potential, then attempt to engineer those trajectories
  \emph{in vivo} or via serial passage.

  \item \textbf{Binding-model exploitation.}
  The calibration map $\Phi_{\mathrm{bind}}$ links mirror-space distances to binding energies.
  Combined with generative models, this could enable direct optimization of antigenic sequences
  to achieve specific binding profiles against known therapeutic antibodies, accelerating
  resistance engineering.
\end{enumerate}

\paragraph{Key distinction from HERALD.}
HERALD detects existing or emergent variants and alerts when dominance risk exceeds a threshold;
it does not optimize over unobserved sequence space. MIRADOR, by contrast, searches over
receptor space to maximize coverage, which implicitly encodes information about which antigenic
states are ``hardest to cover.'' This asymmetry is the source of dual-use risk.

\subsection{Information-hazard mitigations aligned with HERALD}
\label{subsec:mitigations}

MIRADOR adopts the same coarse-grained reporting and access-control principles established for
HERALD, adapted to the design setting.

\paragraph{Coarse-grained outputs.}
\begin{itemize}
  \item \textbf{Panel-level certificates, not singleton solutions.}
  MIRADOR reports coverage certificates
  $\underline{C}_{\mathrm{neut}}(V;L_A^\star)$ for \emph{panels} $V$ of size $K \ge 3$, not for
  individual ``magic bullet'' receptors. This hedges over drift uncertainty and obscures which
  single receptor would be optimal for a specific future variant.

  \item \textbf{Epitope-region summaries, not mutation recipes.}
  Path families and drift forecasts are summarized by broad antigenic regions (e.g., ``RBD
  class 1/2/3 escape,'' ``NTD-focused drift'') rather than combinatorial mutation lists. This
  preserves the scientific value of drift modeling while withholding synthesis-ready sequences.

  \item \textbf{No explicit optimal-antigen disclosure.}
  When MIRADOR identifies a high-coverage panel $V^\star$, the system \emph{does not} publish
  the antigen or path $a_{\mathrm{worst}}$ that is least covered by $V^\star$. Any residual
  gap analysis is reported in aggregate (e.g., ``coverage drops below $60\%$ on paths longer
  than $L_{\max}$'') without sequence-level detail.
\end{itemize}

\paragraph{Restricted access to full geometry.}
The complete antigen--receptor co-embedding $(u_A, u_R)$, including encoder weights and the
binding calibration $\Phi_{\mathrm{bind}}$, is treated as a controlled asset:
\begin{itemize}
  \item \emph{Inference-only APIs} provide coverage certificates for user-supplied panels
        without exposing the embedding functions.
  \item \emph{Audited research access} to the full geometry is restricted to vetted groups
        under data-use agreements that prohibit optimization over hypothetical antigens.
  \item \emph{Governance oversight} analogous to HERALD's biosecurity review applies to any
        release of trained models or path-family datasets.
\end{itemize}

\subsection{Banned optimization modes}
\label{subsec:banned-modes}

To enforce the principle that MIRADOR assists defense against \emph{known or plausibly
forecasted} threats but does not enable \emph{de novo} threat engineering, the following
optimization modes are explicitly prohibited in any deployment:

\begin{enumerate}
  \item \textbf{Unconstrained worst-case antigen search.}
  The system must never accept queries of the form:
  \[
    \text{``Given panel } V, \text{ find the antigen } a_{\mathrm{evade}}
    \text{ that maximally evades } V \text{ under the mirror geometry.''}
  \]
  Such queries directly output immune-evading sequences and are \emph{categorically banned}.

  \item \textbf{Receptor optimization for unobserved variants without input constraints.}
  MIRADOR may optimize receptors $r \in \mathcal{M}_R$ to cover a \emph{specified} set of
  antigens $\mathcal{A}_{\mathrm{focus}}$ provided by the user. It must not accept
  unconstrained prompts like:
  \[
    \text{``Design the optimal antibody for future SARS-CoV-2 variants''}
  \]
  without an explicit path distribution $\Pi_A(L_A^\star)$ derived from surveillance data.
  The requirement is \emph{input-tied optimization}: users supply antigens or path families;
  MIRADOR evaluates or optimizes receptors conditional on those inputs, but does not
  autonomously generate hypothetical future threats.

  \item \textbf{Generative extrapolation far outside benign path families.}
  Generative models may explore within pre-audited benign path families
  $\mathcal{P}_A(L_A^\star)$, but must not extrapolate into regimes that combine:
  \begin{itemize}
    \item large antigenic distances from any observed variant,
    \item high predicted infectivity or virulence (if such metadata exist), and
    \item low predicted coverage by existing or candidate panels.
  \end{itemize}
  In practice, this is enforced via \emph{path-family gates}: any proposed antigen path must
  pass distortion and feasibility checks before being used in coverage optimization.

  \item \textbf{Inverse MIRADOR: receptor-to-antigen generative queries.}
  Tools that condition generative models on a desired receptor-binding profile and output
  corresponding antigenic sequences are \emph{forbidden}. This includes:
  \begin{itemize}
    \item ``Generate an antigen that binds receptor $r$ with affinity $K_D \approx X$.''
    \item ``Find the nearest antigen in sequence space to a given $z_A \in \mathcal{M}_A$.''
  \end{itemize}
  Such capabilities are dual-use by design and must not be exposed in any public or
  semi-public interface.
\end{enumerate}

\subsection{Dual-use firewall in the design loop}
\label{subsec:design-firewall}

Beyond banning specific queries, MIRADOR enforces a \emph{dual-use firewall} at the
architectural level, ensuring that the optimization loop itself cannot be co-opted for threat
generation.

\paragraph{Input-tied optimization only.}
All receptor-panel searches proceed \emph{conditional} on antigens or path distributions
provided as input:
\begin{equation}
  V^\star \in \arg\max_{V \in \mathcal{V}}
              \underline{C}_{\mathrm{neut}}\bigl(V; \Pi_A(L_A^\star)\bigr),
\end{equation}
where $\Pi_A(L_A^\star)$ is either:
\begin{itemize}
  \item derived from HERALD or another surveillance system with documented provenance, or
  \item explicitly specified by the user from a curated library of historical or validated
        drift scenarios.
\end{itemize}
MIRADOR does \emph{not} autonomously hypothesize $\Pi_A(L_A^\star)$ in the absence of data.

\paragraph{Panel-level hedging enforced by design.}
The objective function and constraints in~\eqref{eq:mirador-vaccine-design} and
\eqref{eq:design-problem} penalize singleton solutions:
\begin{itemize}
  \item Minimum panel size $K_{\min} \ge 3$ is enforced.
  \item Diversity constraints (e.g., minimum pairwise mirror-space distance among panel
        members, or epitope-class requirements) prevent collapse to a single ``super-binder.''
  \item Coverage certificates are computed at the \emph{panel} level; individual receptor
        performance is reported only in aggregate.
\end{itemize}
This architectural choice ensures that MIRADOR outputs are inherently hedged and do not reveal
a unique vulnerability axis in antigen space.

\paragraph{Restricted generative sampling.}
When generative models are used for receptor design (Section~\ref{subsec:antibody-design}) or
for augmenting antigen path families, the following safeguards apply:
\begin{enumerate}
  \item \textbf{In-distribution enforcement.}
  Generated sequences must satisfy $a \in \mathcal{I}^A$ or $r \in \mathcal{I}^R$ where
  $\mathcal{I}^A$ and $\mathcal{I}^R$ are pre-audited in-distribution regions. Sequences
  outside these regions trigger OOD gates and are rejected.

  \item \textbf{Benign-path constraint.}
  Generated antigen paths must lie in $\mathcal{P}_A(L_A^\star)$ and satisfy distortion bounds
  $\varepsilon_A(L_A^\star) < \varepsilon_A^{\max}$. Paths that exceed validated distortion are
  flagged and excluded from coverage computations.

  \item \textbf{No unrestricted ``dream'' sequences.}
  Generative models are not permitted to optimize antigen sequences \emph{de novo} against
  arbitrary objectives (e.g., ``maximize immune escape and transmissibility''). All generation
  is conditioned on realistic drift constraints or user-provided anchor variants.
\end{enumerate}

\paragraph{No inverse-query interface.}
The system architecture prohibits any API endpoint or user interface that accepts a receptor
(or panel) as input and returns an antigen (or path) that evades it. Operationally:
\begin{itemize}
  \item Forward queries (``evaluate panel $V$ against path family $\Pi_A$'') are permitted.
  \item Inverse queries (``given panel $V$, find worst-case antigen'') are \emph{hard-blocked}
        at the interface level and flagged for security review if attempted.
\end{itemize}

\subsection{Governance and the MIRADOR dossier}
\label{subsec:governance-dossier}

Every MIRADOR deployment must maintain a \emph{MIRADOR dossier} that documents not only the
geometry and coverage certificates (Section~\ref{sec:discussion}) but also the dual-use
mitigations and oversight structures. At minimum, the dossier includes:

\begin{enumerate}
  \item \textbf{Dual-use risk assessment.}
  A documented evaluation of how the specific path families, antigen library, and receptor
  design space used in the deployment could inform adversarial misuse, together with the
  mitigations applied (coarse-grained reporting, access control, banned modes).

  \item \textbf{Access control and audit logs.}
  A record of who has access to the full mirror geometry $(u_A, u_R, \Phi_{\mathrm{bind}})$,
  under what data-use agreements, and with what oversight. Audit logs track all queries to the
  system, especially any attempts at inverse optimization or out-of-distribution generation.

  \item \textbf{Governance review.}
  Analogous to HERALD's biosecurity oversight, MIRADOR requires review by an institutional
  biosafety committee or equivalent body before:
  \begin{itemize}
    \item releasing trained models or path-family datasets,
    \item publishing detailed coverage certificates that include sequence-level antigen
          information, or
    \item deploying generative capabilities in new pathogen domains.
  \end{itemize}

  \item \textbf{Incident response and misuse monitoring.}
  Procedures for detecting and responding to misuse (e.g., attempts to reverse-engineer
  worst-case antigens from published panels, unauthorized use of the embedding for threat
  modeling). This includes monitoring for papers or preprints that cite MIRADOR in a dual-use
  context and engaging with authors or journals as appropriate.

  \item \textbf{Update and re-review cadence.}
  As surveillance data evolve and new path families are incorporated, the dossier is updated
  and the dual-use risk assessment is re-run. Major updates trigger governance re-review.
\end{enumerate}

The goal is not to prevent MIRADOR's use in legitimate vaccine and therapeutic design, but to
ensure that every deployment operates under explicit dual-use constraints and with documented
accountability.

\subsection{Limitations and open directions}
\label{subsec:ethics-limitations}

Several challenges remain open or only partially addressed by the current framework.

\paragraph{Limits of complementary embeddings.}
The complementary co-embedding assumes a relatively simple binding model (pairwise
antigen--receptor distances). In reality:
\begin{itemize}
  \item \emph{Multi-epitope binding} (avidity, cooperativity, epitope masking) is not fully
        captured by pairwise distances in $\mathcal{M}_A \times \mathcal{M}_R$.
  \item \emph{Polyspecific antibodies} that bind multiple unrelated antigens can confound the
        lock-and-key geometry, especially if the training data include cross-reactive binders
        without adequate annotation.
\end{itemize}
When the embedding breaks down, coverage certificates may be vacuous, but they may also
inadvertently encode spurious vulnerabilities. Better models and slice-aware reporting can
mitigate this, but the risk remains.

\paragraph{Heavy-tailed assay noise and efficacy uncertainty.}
Binding and neutralization assays have fat-tailed error distributions, especially for weak or
intermediate binders. The error budget $E_{\mathrm{bind}}$ and $E_{\mathrm{link}}$ average over
this noise, but extreme outliers (e.g., a highly potent binder that fails to neutralize due to
rare conformational escape) may not be well-characterized. This can lead to overconfident
certificates or, conversely, to panels that appear suboptimal when they are actually robust.

\paragraph{Learning path families vs. hand-coding them.}
Current MIRADOR deployments assume that $\mathcal{P}_A(L_A^\star)$ and
$\mathcal{P}_R(L_R^\star)$ are specified by domain experts or inherited from surveillance
systems. Learning path families directly from data---subject to interpretability and dual-use
constraints---is an active research direction. Unsupervised or weakly supervised path learning
could improve forecast quality but also risks discovering ``natural-looking'' high-risk
trajectories that should remain undisclosed.

\paragraph{Interactions with human decision-making.}
Coverage certificates are probabilistic lower bounds, not guarantees. In practice, decision-makers
may:
\begin{itemize}
  \item over-interpret certificates as hard assurances, leading to inadequate preparedness when
        paths fall outside the validated regime;
  \item under-utilize certificates due to mistrust of model-based reasoning, defaulting to
        reactive strain-matching; or
  \item use certificates in adversarial ways (e.g., optimizing for regulatory approval rather
        than public-health impact).
\end{itemize}
Understanding and shaping these human-AI interactions is critical for MIRADOR's responsible
deployment and remains a largely open sociotechnical problem.

\paragraph{Evolving dual-use landscape.}
As generative biology capabilities advance, the threat model for MIRADOR will shift. What is a
banned optimization mode today may become a standard tool tomorrow, requiring ongoing re-evaluation
of firewalls and access policies. MIRADOR's governance framework must be adaptive, with regular
threat-model updates and community input from biosecurity, public health, and AI safety domains.

\subsection{Summary of ethical commitments}

MIRADOR's ethical stance is:
\begin{enumerate}
  \item \textbf{Design for defense, not offense.}
  All optimizations are input-tied and conditional on user-supplied antigens or surveillance-derived
  path families. The system does not autonomously generate or optimize hypothetical threats.

  \item \textbf{Panel-level outputs, not singleton solutions.}
  Coverage certificates and receptor designs are always reported for \emph{panels} with enforced
  diversity, obscuring singular vulnerability axes.

  \item \textbf{Coarse-grained public reporting.}
  Detailed sequence-level or path-level information is restricted; public outputs summarize
  coverage at the epitope-region or lineage level.

  \item \textbf{Hard-blocked inverse queries.}
  Receptor-to-antigen inversion and worst-case evading-antigen searches are architecturally
  prohibited and flagged for review if attempted.

  \item \textbf{Governance and accountability.}
  Every deployment maintains a MIRADOR dossier with dual-use risk assessment, access control,
  oversight review, and incident-response procedures.

  \item \textbf{Adaptive re-evaluation.}
  As capabilities and threats evolve, the framework's firewalls and policies are periodically
  re-assessed with input from biosecurity and public-health communities.
\end{enumerate}

These commitments cannot eliminate dual-use risk entirely, but they establish a baseline of
responsible practice and a culture of reflexive scrutiny. MIRADOR is a tool for proactive public
health; its design and governance aim to ensure it remains so.

\section{Discussion and Outlook}
\label{sec:discussion}

This section summarizes how the mirror geometry connects to receptor-panel design, how coverage
certificates should be interpreted in practice, and where the framework is fragile or incomplete.

\subsection{Why mirror geometry matters for design}

The construction in Sections~2--5 is deliberately geometric. For design, three ingredients are
central:

\begin{enumerate}
  \item \emph{Antigen side.}
  The antigen manifold $(\mathcal{M}_A, g_A)$, with benign path family
  $\mathcal{P}_A(L_A^\star)$ and distortion profile $\varepsilon_A(L_A^\star)$, encodes what the
  virus can plausibly do over the operational horizon $L_A^\star$. The forecast path law
  $\Pi_A(L_A^\star)$ specifies \emph{how} it is likely to move within that manifold.

  \item \emph{Receptor side.}
  The receptor manifold $(\mathcal{M}_R, g_R)$, with path family $\mathcal{P}_R(L_R^\star)$ and
  distortion profile $\varepsilon_R(L_R^\star)$, encodes what receptor states are biologically
  achievable and immunologically acceptable over a horizon $L_R^\star$ (e.g., affinity
  maturation or design moves). The in-distribution region $\mathcal{I}^R \subseteq \mathcal{M}_R$
  delimits receptors supported by data and used in coverage certificates.

  \item \emph{Lock--key geometry.}
  The complementary co-embedding (Definition~\ref{def:complementary-coembedding}) ties the two
  manifolds together: for $a \in \mathcal{M}_A$ and $r \in \mathcal{M}_R$, the ambient distance
  $\delta(a,r)$ is a calibrated surrogate for binding energy and, via $E_{\text{link}}$, for
  neutralization. Pathwise hit events and coverage (Definitions~\ref{def:mirror-paths} and
  \ref{def:hit-event}) then become geometric statements about when antigen paths in
  $\mathcal{P}_A(L_A^\star)$ come within a binding radius of receptors in a panel $V \subseteq
  \mathcal{M}_R$.
\end{enumerate}

HERALD and similar Davis-style surveillance systems estimate \emph{where} the virus can go and how
it is drifting. Within the same geometric language, MIRADOR answers a different question:
\emph{which receptor panels come with meaningful, slice-aware coverage guarantees against those
drift trajectories, on a pre-registered horizon and in a specified in-distribution region?}

\subsection{Panel design as a constrained coverage problem}

Let $\mathcal{V}$ denote a design space of feasible receptor panels (e.g., all subsets of a candidate
library satisfying manufacturability and diversity constraints). Let
$V \subseteq \mathcal{I}^R$ be a candidate panel whose members lie in the binding in-distribution
region of the receptor manifold. For a fixed antigen path horizon $L_A^\star$, define
\[
  C_{\mathrm{neut}}(V;L_A^\star)
  := \Pr_{\Gamma_A \sim \Pi_A(L_A^\star)}\!\bigl(
       \Hit_{\mathrm{neut}}(V,\Gamma_A)
     \bigr),
\]
the (unknown) neutralization coverage of $V$ at horizon $L_A^\star$ under the path law
$\Pi_A(L_A^\star)$.

Section~5 introduced a \emph{coverage certificate}
\[
  \underline{C}_{\mathrm{neut}}(V;L_A^\star)
  :=
  \bigl[\widehat{C}(V;L_A^\star) - \varepsilon_{\mathrm{samp}}\bigr]_+
  - \bigl(\overline{E}_{\mathrm{geom}} + \overline{E}_{\mathrm{bind}} + \overline{E}_{\mathrm{link}}\bigr),
\]
where $[u]_+ := \max\{u,0\}$, $\widehat{C}(V;L_A^\star)$ is an empirical coverage estimate based on
Monte Carlo paths, $\varepsilon_{\mathrm{samp}}$ is a sampling tolerance chosen from a concentration
inequality, and $\overline{E}_{\mathrm{geom}},\overline{E}_{\mathrm{bind}},\overline{E}_{\mathrm{link}}$ are
pre-registered upper bounds on the geometric, binding, and linkage error components. Under the
assumptions of Theorem~5.1, this certificate is a high-probability \emph{one-sided} lower bound:
\[
  C_{\mathrm{neut}}(V;L_A^\star)
  \;\ge\;
  \underline{C}_{\mathrm{neut}}(V;L_A^\star)
  \quad\text{with high probability.}
\]

Given a policy-relevant target $C_\star \in (0,1)$ (e.g., ``at least $70\%$ neutralization coverage
against forecast drift paths''), the design problem can be written as
\begin{equation}
\label{eq:design-problem}
  \text{find } V \in \mathcal{V}
  \quad\text{such that}\quad
  \underline{C}_{\mathrm{neut}}(V;L_A^\star) \;\ge\; C_\star.
\end{equation}
In words: choose a panel whose certified coverage, after subtracting sampling and model-error
budgets, still exceeds the target.

This framing deliberately separates three roles:
\begin{enumerate}
  \item $C_{\mathrm{neut}}(V;L_A^\star)$ is the unobservable ground-truth coverage under the path
  law.
  \item $\widehat{C}(V;L_A^\star)$ is a Monte Carlo estimate, with sampling uncertainty controlled
  by $\varepsilon_{\mathrm{samp}}$.
  \item $\underline{C}_{\mathrm{neut}}(V;L_A^\star)$ is the conservative, pre-registered quantity
  used in design and governance.
\end{enumerate}
The certificate is intentionally conservative and may be vacuous when error budgets or independence
slack are large; the framework treats such vacuity as a reason to abstain from strong claims, not as
a parameter to be tuned away.

\subsection{A practical workflow for MIRADOR-style design}

A MIRADOR-style deployment can be organized as the following high-level workflow.

\begin{enumerate}
  \item \textbf{Fix scope and candidate pool.}
  Specify the application (e.g., prophylactic antibody panel, vaccine-adjacent receptor library),
  the time horizon $L_A^\star$ of interest, and the design space $\mathcal{V}$ (e.g., all panels of
  size $m$ drawn from a curated receptor set).

  \item \textbf{Construct or inherit mirror geometry.}
  Train or fine-tune encoders to obtain $(\mathcal{M}_A,g_A)$ and $(\mathcal{M}_R,g_R)$ together
  with the complementary co-embedding (Definition~\ref{def:complementary-coembedding}). Audit
  pathwise distortion and configuration margins to verify non-vacuity on candidate horizons
  $L_A^\star, L_R^\star$.

  \item \textbf{Identify the binding in-distribution region.}
  Use training coverage, held-out binding data, and geometry-based quality checks to define a
  binding in-distribution region
  $\mathcal{I}^A \times \mathcal{I}^R \subseteq \mathcal{M}_A \times \mathcal{M}_R$. Define quality
  and out-of-distribution (OOD) gates that enforce:
  \[
    (a,r) \text{ used for coverage certificates}
    \;\Longrightarrow\;
    (a,r) \in \mathcal{I}^A \times \mathcal{I}^R.
  \]

  \item \textbf{Choose path families and operational horizons.}
  Specify benign path families
  $\mathcal{P}_A(L_A^\star) \subseteq C^1\bigl([0,1],\mathcal{M}_A\bigr)$ and
  $\mathcal{P}_R(L_R^\star) \subseteq C^1\bigl([0,1],\mathcal{M}_R\bigr)$ that reflect realistic
  antigenic drift and receptor evolution. Perform distortion audits to ensure that
  $\varepsilon_A(L_A^\star)$ and $\varepsilon_R(L_R^\star)$ are small enough that the non-vacuity
  conditions remain satisfied.

  \item \textbf{Estimate the antigen path law.}
  Using HERALD or another Davis-style surveillance system, estimate the path distribution
  $\Pi_A(L_A^\star)$ restricted to $\mathcal{I}^A$ and the forecast slice of interest (e.g., lineage,
  geography, or time window). Document how $\Pi_A(L_A^\star)$ was fit and how forecast
  uncertainty is handled.

  \item \textbf{Search over panels and compute coverage certificates.}
  For $V \in \mathcal{V}$, draw Monte Carlo paths $\Gamma_A^{(b)} \sim \Pi_A(L_A^\star)$, compute
  the empirical hit fraction $\widehat{C}(V;L_A^\star)$, and evaluate the certificate
  $\underline{C}_{\mathrm{neut}}(V;L_A^\star)$. This can be done either via heuristic search (e.g.,
  greedy or evolutionary algorithms) or over a restricted design grid. Slice-aware certificates
  (e.g., per-lineage or per-region) should be computed in parallel.

  \item \textbf{Select and document.}
  Choose one or more panels $V^\star$ that satisfy~\eqref{eq:design-problem} on the intended
  slices. For each panel, record the geometric assumptions, in-distribution region,
  $\Pi_A(L_A^\star)$, error-budget estimates, and coverage certificates that justify the choice.
\end{enumerate}

The workflow intentionally separates \emph{what} is being guaranteed (pathwise coverage against
$\Pi_A(L_A^\star)$ on $\mathcal{I}^A \times \mathcal{I}^R$) from \emph{how} those guarantees are
estimated (Monte Carlo paths, distortion audits, error-budget estimation).

\subsection{Integration with HERALD and other surveillance systems}

The mirror geometry is most useful when coupled to a surveillance system that continuously updates
$\Pi_A(L_A^\star)$ and monitors when the antigen process leaves its validated regime.

HERALD provides exactly this: a Davis-style antigen manifold, a drift statistic, and a compositional
error budget for dominance-risk alerts. In a joint deployment, one can envisage the following
closed loop:
\begin{enumerate}
  \item HERALD runs continuously and maintains a forecast path law $\Pi_A(L_A^\star)$ over
  $(\mathcal{M}_A,g_A)$, together with slice-aware error budgets for geometry, linkage, and
  calibration.

  \item MIRADOR uses $\Pi_A(L_A^\star)$ and its uncertainty envelope as input, together with a
  receptor manifold $(\mathcal{M}_R,g_R)$ and in-distribution region $\mathcal{I}^R$, to design
  panels with certified coverage on the current forecast.

  \item As new data arrive, HERALD detects when empirical drift paths move outside
  $\mathcal{I}^A \subseteq \mathcal{M}_A$ or when its own error budget becomes vacuous. In those
  regimes, pathwise coverage certificates for MIRADOR panels must be treated as out-of-scope:
  either the geometric assumptions must be revalidated or coverage claims must be narrowed.
\end{enumerate}

Conceptually, HERALD answers ``where might the virus go, on what timescale, and with what
uncertainty?'' while MIRADOR answers ``given that forecast, which receptor panels can plausibly
intercept those paths with non-vacuous, slice-aware coverage certificates?'' The two systems share
the same geometric language (Davis manifolds, path families, error budgets), which makes joint
reasoning and auditing possible.

\subsection{Limitations and failure modes specific to MIRADOR}

In addition to the general limitations of Davis systems, MIRADOR introduces several design-specific
failure modes.

\paragraph{Dependence on the antigen path law.}
Coverage certificates are conditional on $\Pi_A(L_A^\star)$ and on remaining in
$\mathcal{I}^A$. If the true antigen process places substantial mass on paths outside the estimated
law or outside $\mathcal{I}^A$ (e.g., due to recombination, strong immune imprinting, or unexpected
host shifts), the true coverage $C_{\mathrm{neut}}(V;L_A^\star)$ may fall below the certificate.
Robustness to misspecified path laws remains an open problem.

\paragraph{Search over a restricted design space.}
The design space $\mathcal{V}$ is constrained by manufacturability, diversity, and biological
plausibility. A certificate for $V^\star \in \mathcal{V}$ says nothing about panels outside
$\mathcal{V}$, and non-existence of a good panel within $\mathcal{V}$ does not imply that no good
panel exists in principle.

\paragraph{Binding and neutralization gaps.}
The linkage term $E_{\mathrm{link}}$ explicitly acknowledges that binding in the co-embedding need
not imply neutralization. When the biology of effector function, Fc engineering, or compartmental
effects is poorly captured, $E_{\mathrm{link}}$ will be large and certificates may become vacuous
even when geometric coverage looks good.

\paragraph{Joint error budgets and independence slack.}
Error components for geometry, binding, linkage, and sampling may be strongly correlated on certain
slices (e.g., particular lineages or capture conditions). In such regimes the independence slack
$\delta_{\mathrm{indep}}^{\mathrm{cov}}$ may dominate, forcing reliance on conservative union bounds and
leading to vacuous certificates. The framework treats this correctly as a signal that stronger claims
should not be made from model outputs alone.

\subsection{Research directions}

Several directions emerge naturally from the coverage-certificate perspective.

\begin{itemize}
  \item \textbf{Robust coverage under path-law uncertainty.}
  Instead of conditioning on a single $\Pi_A(L_A^\star)$, one could seek certificates that hold
  uniformly over a family of path laws constrained by surveillance data, or that explicitly fold
  path-law uncertainty into the error budget.

  \item \textbf{Learning or adapting path families.}
  In the present work, $\mathcal{P}_A(L_A^\star)$ and $\mathcal{P}_R(L_R^\star)$ are designed
  objects. Learning path families subject to curvature and plausibility constraints---or adapting
  them as data accumulate---could reduce misspecification while retaining interpretability.

  \item \textbf{Multi-epitope and multi-compartment coverage.}
  Real-world protection often depends on coordinated coverage across epitopes, compartments, or
  modalities. Extending coverage certificates to multi-layer panels (e.g., systemic + mucosal, or
  antibody + T-cell) is mathematically straightforward but practically challenging.

  \item \textbf{Search procedures with guarantees.}
  Design algorithms that explore $\mathcal{V}$ while preserving monotonicity of coverage
  certificates, or that provide \emph{search} guarantees (e.g., approximate optimality with respect
  to~\eqref{eq:design-problem}), would strengthen the connection between the theory and concrete
  panel proposals.

  \item \textbf{Benchmarks and community standards.}
  Having multiple groups compute coverage certificates for shared geometry + path-law benchmarks
  would sharpen understanding of when MIRADOR-style designs are robust and when they are fragile.
\end{itemize}

\subsection{Governance and the MIRADOR dossier}

Coverage certificates are only as trustworthy as the documentation around them. For any deployed
MIRADOR-style system, we recommend maintaining a \emph{MIRADOR dossier} that includes, at
minimum:

\begin{enumerate}
  \item \textbf{Construction details for $(\mathcal{M}_A, g_A)$ and $(\mathcal{M}_R, g_R)$.}
  Encoder architectures, training objectives, regularization schemes, and distortion audits that
  justify the bounded-distortion regime on $\mathcal{P}_A(L_A^\star)$ and $\mathcal{P}_R(L_R^\star)$.

  \item \textbf{In-distribution regions and horizons.}
  Formal definitions of $\mathcal{I}^A$ and $\mathcal{I}^R$, the criteria used to enforce
  in-distribution constraints at training and deployment time, and the chosen horizons
  $L_A^\star, L_R^\star$ with supporting distortion and non-vacuity evidence.

  \item \textbf{Binding model and calibration.}
  The data and procedures used to fit the binding and neutralization mappings (including
  $\Phi_{\mathrm{bind}}$ and $\Delta G_{\mathrm{thresh}}$), together with validation diagnostics and slice-wise
  performance.

  \item \textbf{Antigen forecast and sampling.}
  A description of how $\Pi_A(L_A^\star)$ is estimated (e.g., via HERALD), including uncertainty
  quantification, slice definitions, and any assumptions about stationarity or mixing.

  \item \textbf{Coverage certificates and design decisions.}
  For each proposed panel $V$, the numerical value of
  $\underline{C}_{\mathrm{neut}}(V;L_A^\star)$ and its components
  $\widehat{C},\varepsilon_{\mathrm{samp}},\overline{E}_{\mathrm{geom}},\overline{E}_{\mathrm{bind}},\overline{E}_{\mathrm{link}},
  \delta_{\mathrm{indep}}^{\mathrm{cov}}$, reported globally and on key slices. Design choices
  (such as panel size, diversity constraints, and excluded candidates) should be tied explicitly
  to these numbers.

  \item \textbf{Monitoring and update policy.}
  Rules for when to recompute certificates (e.g., after substantial changes in $\Pi_A$ or in
  error-budget estimates), criteria for declaring certificates vacuous on certain slices, and the
  fallback procedures (laboratory assays, human review) that apply in those regimes.

  \item \textbf{Limitations and disclaimers.}
  A clear statement of assumptions, known blind spots (e.g., recombination, novel antigenic
  mechanisms), and the intended decision context, so that end users do not over-interpret what a
  MIRADOR certificate guarantees.
\end{enumerate}

The goal is not to treat MIRADOR outputs as oracles, but to make explicit \emph{which} geometric and
statistical assumptions underwrite each coverage claim, and to provide the artifacts needed for
external review. In that sense, coverage certificates are not just numbers but contracts: they bind
the system to a specific Davis-style regime of validity, with clearly delineated conditions under
which abstention, revision, or re-design is required.


\end{document}